\usepackage{../submodules/FancyBeamer/fancybeamer} % use the fancy beamer package
%\fancylogos{,uulm,,unibe,,ovgu-blue,}
\setfancycolumnsdefault{animation=keep} % animate all columns per default

% GRAPHICSPATH

\setpaths{
	{../pics/}
	{../submodules/Title-Pictures/}
	{../submodules/Course-on-Software-Product-Lines/pics/}
	% Florian Sihler's additions
	{../submodules/abstract-interpretation-ltx}
	{../submodules/code-animation}
	{../submodules/xlistings/langs}
}

% IMPORTED PACKAGES

\usepackage{multicol} % used temporarily for the lecture overview
\usepackage{stmaryrd} % \lightning in modeling lecture
\usepackage{catchfile} % read picture sources from files
\usepackage{silence} % ignore some warnings
\WarningFilter{latex}{You have requested package} % ignore warning for loading packages from submodules
%\usepackage{../EagleoutIce-fancyqr/fancyqr} % QR codes for practice task, credits to Florian Sihler
\usepackage{fancyqr} % QR codes for practice task, credits to Florian Sihler
\fancyqrset{gradient=true,gradient angle=-45,l color=black,r color=blue}
\usepackage{ifthen} % \ifuniversity
\usepackage{booktabs} % tables
%\usepackage{colortbl} % \rowcolor used in Lecture 8 in SPL
%\usepackage{algpseudocode} % used in Lecture 10 in SPL

% Florian Sihler's crazy additions because he ported his slides and tries to be a special snowflake
\usepackage{../submodules/abstract-interpretation-ltx/absint}
\usepackage{../submodules/code-animation/code-animation}
\usepackage{../submodules/xlistings/xlistings}
\makeatletter\LoadLanguages{java,R}
\usepackage{stmaryrd} % flo's symbols
\xlstsetmintedstyle{plain}
\def\Set#1{\ensuremath{\{#1\kern1pt\}}}
\def\absintstyle#1{#1}
\def\IntCC#1#2{\ensuremath{[#1\,..\,#2]}}
\def\IntOC#1#2{\ensuremath{(#1\,..\,#2]}}
\def\IntCO#1#2{\ensuremath{[#1\,..\,#2)}}
\def\IntOO#1#2{\ensuremath{(#1\,..\,#2)}}
\def\AbstractInfo#1{\ensuremath{\mathcolor{gray}{\Lbag}\,#1\,\mathcolor{gray}{\Rbag}}}
\usepackage{forest}

\usepackage{csquotes}
\usepackage[style=alphabetic,backend=biber]{biblatex}
\usepackage{multirow}
\usepackage{tikzpingus}
\usepackage{relsize}

% INCLUDED TEMPLATE FILES

% TYPICAL COMMANDS FOR LECTURES

\renewcommand{\emph}[1]{{\color{blue}\textbf{#1}}}
%\ifdarkmode
%	\renewcommand{\emph}[1]{{\color{orange}\textbf{#1}}} % workaround for better readability
%\fi

\newcommand{\deutsch}[1]{{\color{gray}(#1)}}
\ifdarkmode
	\renewcommand{\deutsch}[1]{{\color{lightgray}(#1)}}
\fi
\newcommand{\deutschertitel}[1]{{\tiny\deutsch{#1}}}

\newcommand{\mycite}[1]{``#1''}
\newcommand{\mycitebegin}{``}
\newcommand{\myciteend}{''}
\newcommand{\mytitlesource}[1]{{\tiny\normalfont\mbox{[#1]}}}
\newcommand{\mysource}[1]{\ifthenelse{\equal{#1}{}}{}{\phantom{.}~\hfill~\mytitlesource{#1}}}
\newcommand{\mychapter}[1]{, Chapter~#1}
\newcommand{\mysection}[1]{, Section~#1}
\newcommand{\mypage}[1]{, p.~#1}
\newcommand{\mypages}[1]{, pp.~#1}

%\newcommand{\todo}[1]{{\color{red}\textbf{[#1]}}}
%\newcommand{\fodo}[1]{\todo{\footnote{\todo{#1}}}}
%\newcommand{\todots}{\todo{\ldots}}

\newcommand{\textheightoftitle}{.125\textheight}
\newcommand{\textheightwithtitle}{.825\textheight}
\newcommand{\textheightwithouttitle}{.95\textheight}

\newcommand{\lessonslearned}[3]{
	\subsection{Summary}
	\begin{frame}{\insertsection{} -- \myframetitle}
		\begin{fancycolumns}[animation=none]
			\begin{definition}{Lessons Learned}
				\begin{itemize}
					#1
				\end{itemize}
			\end{definition}
			\uncover<2->{\begin{note}{Further Reading}
				\small % references take space, can be a little smaller
				\begin{itemize}
					#2
				\end{itemize}
			\end{note}}
		\nextcolumn
			\uncover<3->{\begin{example}{Practice}
				#3
			\end{example}}
		\end{fancycolumns}
	\end{frame}
}

% TODO would be great to also show the section titles on this slide, but not sure how to do this
\newcommand{\faq}[3]{
	\subsection{FAQ}
	\begin{frame}{\emph{FAQ} -- \insertsubtitle{}}
		\small 
		\begin{fancycolumns}[t,columns=3,widths={33,33,33},animation=none]
			\emph{\large Lecture~\insertlecturenumber\parta}
			\begin{itemize}
				#1
			\end{itemize}
		\nextcolumn
			\emph{\large Lecture~\insertlecturenumber\partb}
			\begin{itemize}
				#2
			\end{itemize}
		\nextcolumn
			\emph{\large Lecture~\insertlecturenumber\partc}
			\begin{itemize}
				#3
			\end{itemize}
		\end{fancycolumns}
	\end{frame}
}

% REFERENCES TO OTHER LECTURES

\newcommand{\emphifequal}[3]{\ifthenelse{\equal{#1}{#2}}{\bfseries}{}\hyperlink{lecture#2}{#3}}

\newcommand{\reflecture}[1]{\hyperlink{lecture#1}{Lecture~#1}}
\newcommand{\parta}{a}
\newcommand{\partb}{b}
\newcommand{\partc}{c}
% TODO how to handle different numberings for different universities?

% COMMANDS FOR PROPOSITIONAL FORMULAS AND MATHEMATICAL NOTATIONS

\newcommand{\sem}[1]{\ensuremath{\llbracket #1 \rrbracket}} % semantics brackets
\newcommand{\pand}{\wedge} % conjunction
\newcommand{\por}{\vee} % disjunction
\newcommand{\pnot}{\neg} % negation
\newcommand{\pequals}{\leftrightarrow} % biconditional
\newcommand{\npequals}{\nleftrightarrow} % exclusive disjunction
\newcommand{\mequals}{\Leftrightarrow} % equivalence (meta-level)
\newcommand{\pimplies}{\rightarrow} % conditional
\newcommand{\mimplies}{\Rightarrow} % implication (meta-level)
\newcommand{\defeq}{\vcentcolon=} % defining equals
\newcommand{\power}[1]{\mathcal{P}(#1)} % power set
\newcommand{\refslide}[1]{\hyperlink{#1}{(see Slide \autoref{#1})}} % link to slide
\newcommand{\fs}[1]{\ensuremath{{\color{green}#1}}} % selected feature
\newcommand{\fd}[1]{\ensuremath{{\color{red}#1}}} % deselected feature
\newcommand{\clause}[1]{\ensuremath{{\color{red}#1}}} % clause in a CNF
\newcommand{\cfg}[3][0]{\ensuremath({{\only<#1>{\color{green}}\{#2\}}, {\only<#1>{\color{red}}\{#3\}})}} % configuration

\usepackage{mathtools} % required for absolute value in modeling lecture
\DeclarePairedDelimiter\abs{\lvert}{\rvert} % absolute value

% COMMANDS TO INCLUDE XKCDs AND PICS

\newcommand{\xkcd}[2]{\picDark[#2]{xkcd/#1}}
\newcommand{\xkcdframe}[1]{
	\begin{frame}
		\centering%
		\xkcd{#1}{height=80mm}
	\end{frame}
}
\newcommand{\widexkcdframe}[1]{
	\begin{frame}
		\xkcd{#1}{width=\linewidth}
	\end{frame}
}

\newcommand{\pictureframe}[2]{
	{
		\usebackgroundtemplate{#1}
		\begin{frame}[plain]
			#2
		\end{frame}
	}
}

% COMMANDS TO SIMPLIFY USAGE OF THE PROJECT CARTOON

\newcommand{\projectcartoonwidth}{.18} % default value for 5 tiles
\newcommand{\projectcartoon}[2]{\projectcartoonimage{sepia/cell_#1}{#2}}
\newcommand{\hprojectcartoon}[2]{\projectcartoonimage{alpha/cell_#1}{#2}}
\newcommand{\projectcartoonimage}[2]{
	\begin{minipage}[t]{\projectcartoonwidth\linewidth}
		\centering
		\pic[width=\linewidth]{projectcartoon/#1}\\#2
	\end{minipage}\hfill
}
\newcommand{\waterfallcartoon}{
	\hprojectcartoon{02}{Analysis}%
	\hprojectcartoon{03}{Design}%
	\hprojectcartoon{04}{Implementation}%
	\hprojectcartoon{05}{Testing}%
	\hprojectcartoon{10}{Maintenance}%
}
%\uncover<1->{\hprojectcartoon{01}{how the customer explained it}} % requirements
%\uncover<1->{\hprojectcartoon{02}{how the project leader understood it}} % modeling
%\uncover<1->{\hprojectcartoon{03}{how the analyst designed it}} % architecture and design
%\uncover<1->{\hprojectcartoon{04}{how the programmer implemented it}} % implementation
%\uncover<1->{\hprojectcartoon{05}{what the beta testers received}} % testing
%\uncover<1->{\hprojectcartoon{06}{how the business consultant described it}} % process
%\uncover<1->{\hprojectcartoon{07}{how the project was documented}} % documentation?
%\uncover<1->{\hprojectcartoon{08}{what operations installed}} % devops/continuous integration?
%\uncover<1->{\hprojectcartoon{09}{how the customer was billed}} % management (pricing)
%\uncover<1->{\hprojectcartoon{10}{how it was supported}} % maintenance
%\uncover<1->{\hprojectcartoon{11}{what marketing advertised}} % product lines. marketing? reuse?
%\uncover<1->{\hprojectcartoon{12}{when it was delivered}} % configuration management. continous delivery?
%\uncover<1->{\hprojectcartoon{13}{what the customer really needed}} % customer / SE II
%\uncover<1->{\hprojectcartoon{14}{what the digg effect can do to your site}} % ???
%\uncover<1->{\hprojectcartoon{15}{the disaster recover plan}} % ???
%\uncover<1->{\hprojectcartoon{16}{the open source version}} % open source (licensing)
%\uncover<1->{\hprojectcartoon{17}{how it performed under load}} % compilation and static analyses. quality assurance? performance?
%\uncover<1->{\hprojectcartoon{18}{how patches were applied}} % evolution

% COMMANDS BASED ON TIKZ

\newcommand{\mglass}[1]{
	\scalebox{#1}{
		\begin{tikzpicture}[overlay]
			\node[circle,draw=gray,line width=1,fill=blue,fill opacity=.2,inner sep=2mm] {} edge[draw=gray,line width=1,double distance=1pt,shorten <= -.2pt] (.33,-.33);
		\end{tikzpicture}
	}
}

\def\checkmark{\tikz\fill[scale=0.4](0,.35) -- (.25,0) -- (1,.7) -- (.25,.15) -- cycle;}

\newcommand{\correct}{%
	\begin{tikzpicture}[yscale=.1,xscale=.1,xshift=20,yshift=-10,overlay]
		\draw[line width=3,line join=bevel] (0,1) to [bend left=30] (1,0) to [bend left=15] (3,3);
		\draw[line width=2,line join=bevel,green] (0,1) to [bend left=30] (1,0) to [bend left=15] (3,3);
	\end{tikzpicture}%
}
\newcommand{\alsocorrect}{%
	\begin{tikzpicture}[yscale=.1,xscale=.1,xshift=20,yshift=-10,overlay]
		\draw[line width=3,line join=bevel] (0,1) to [bend left=30] (1,0) to [bend left=15] (3,3);
		\draw[line width=2,line join=bevel,orange] (0,1) to [bend left=30] (1,0) to [bend left=15] (3,3);
	\end{tikzpicture}%
}
\newcommand{\wrong}{%
	\begin{tikzpicture}[yscale=.1,xscale=.1,xshift=20,yshift=-10,overlay]
		\draw[line width=3,line join=bevel] (0,2) to [bend left=5] (2,0) (0,0) to [bend left=15] (2,2);
		\draw[line width=2,line join=bevel,red] (0,2) to [bend left=5] (2,0) (0,0) to [bend left=15] (2,2);
	\end{tikzpicture}%
}

% reusable tikz pictures
\tikzset{
	versioncontrol/.style={xscale=.55,font=\footnotesize,auto,knoten/.style={circle,inner sep=0,outer sep=0},
		kante/.style={overlay,->,thick,bend right=5,shorten >=-5pt,shorten <=-5pt},
		innerekante/.style={overlay,->,thick,bend right=60,shorten >=-10pt,shorten <=-10pt}}
}
\newcommand{\clientserver}[7]{
	\begin{center}
		\begin{tikzpicture}[versioncontrol]
			\node[knoten] (server) at (0,0) {\picDark[scale=.15]{versioncontrol/#1} \pic[scale=.3]{versioncontrol/tonne}};
			\node[knoten] (client1) at (-3,-2) {~~~~~\picDark[scale=.1]{versioncontrol/#2}~~~~~};
			\node[knoten] (client2) at (3,-2) {~~~~~\picDark[scale=.1]{versioncontrol/#3}~~~~~};
			#7
			\node[overlay] at (0,.8) {#4};
			\node[overlay] at (-3,-2.65) {#5};
			\node[overlay] at (3,-2.65) {#6};
		\end{tikzpicture}
	\end{center}
}

%scale 0.15 = 30pt
\newcommand{\peertopeer}[8]{
	\vspace{-2mm}
	\begin{center}
		\begin{tikzpicture}[versioncontrol]
			\node[knoten] (server) at (0,0) {\picDark[height=30pt]{versioncontrol/#1} \pic[height=20pt]{versioncontrol/#2}};
			\node[knoten] (client1) at (-3,-2) {\picDark[height=20pt]{versioncontrol/#3} \pic[height=15pt]{versioncontrol/#2}};
			\node[knoten] (client2) at (3,-2) {\picDark[height=20pt]{versioncontrol/#4} \pic[height=15pt]{versioncontrol/#2}};
			\node[knoten] (client21) at (3.6,-1.75) {};
			#8
			\node[above = -12pt of server, overlay] {#5};
			\node[text width = 0.45\linewidth, align = center, below = 2pt of client1, overlay] {#6};
			\node[text width = 0.45\linewidth, align = center, below = 2pt of client2, overlay] {#7};
		\end{tikzpicture}
	\end{center}
	\vspace{5mm}
}

\newcommand{\trunkbranch}[2]{
	\begin{tikzpicture}[auto,xscale=.75,font=\footnotesize,trunk/.style={circle,draw,thick,minimum size=2mm,inner sep=0mm,outer sep=.5mm},branch1/.style={trunk,fill=blue,yshift=1cm},branch2/.style={trunk,fill=blue,yshift=2cm},change/.style={overlay,->,thick},main/.style={trunk}, branch/.style={change#1},merge/.style={change#2}]
		\node[main,label=below:r1,label=left:main] at (0,0) (rev1) {};
		\node[branch1,label=above:r2] at (1,0) (rev2) {};
		\node[branch1,label=above:r3] at (2,0) (rev3) {};
		\node[main,label=below:r4] at (3,0) (rev4) {};
		\node[branch2,label=above:r5,label=left:a branch] at (4,0) (rev5) {};
		\node[branch1,label=above:r6] at (5,0) (rev6) {};
		\node[branch1,label=above:r7] at (6,0) (rev7) {};
		\node[branch2,label=above:r8] at (7,0) (rev8) {};
		\node[main,label=below:r9] at (8,0) (rev9) {};
		\node[branch2,label=above:r10] at (9,0) (rev10) {};
		\node[main,label=below:r11] at (10,0) (rev11) {};
		\node[main,label=below:r12] at (11,0) (rev12) {};
		\node[branch2,label=above:r13] at (12,0) (rev13) {};
		\draw[branch] (rev1) to (rev2);
		\draw[change] (rev1) to (rev4);
		\draw[change] (rev2) to (rev3);
		\draw[merge] (rev3) to (rev4);
		\draw[branch] (rev4) to (rev5);
		\draw[branch] (rev4) to (rev6);
		\draw[change] (rev4) to (rev9);
		\draw[change] (rev5) to (rev8);
		\draw[change] (rev6) to (rev7);
		\draw[merge] (rev7) to (rev9);
		\draw[change] (rev8) to (rev10);
		\draw[merge] (rev9) to (rev10);
		\draw[change] (rev9) to (rev11);
		\draw[merge] (rev10) to node {merge} (rev11);
		\draw[change] (rev10) to (rev13);
		\draw[change] (rev11) to (rev12);
	\end{tikzpicture}
}

% OTHER COMMANDS

\usetikzlibrary{snakes}
\newcommand{\tpkt}[1]{\tikz[remember picture] \coordinate(#1);}
\newcommand{\myunderline}[2]{
	\begin{tikzpicture}[remember picture,overlay]
		\foreach \x in {#1} {
			\node[yshift=.5mm] (x1) at (\x 1) {};
			\node[yshift=.5mm] (x2) at (\x 2) {};
			\draw[#2,thick,snake=snake,segment amplitude=.2mm,segment length=1mm] (x1.south) -- (x2.south);
		}
	\end{tikzpicture}
}
\newcommand{\errorunderline}[1]{\myunderline{#1}{red}}
\newcommand{\warningunderline}[1]{\myunderline{#1}{orange}}
\newcommand{\infounderline}[1]{\myunderline{#1}{blue}}

\usepackage{listings}
\usepackage[normalem]{ulem}
\tcbuselibrary{listings}

\newcommand{\definelanguage}[2]{
	\lstdefinelanguage{#1}{#2}
	\lstnewenvironment{#1}[1][]{\lstset{language=#1,##1}}{}
	\newtcblisting{#1tight}[2][]{fancy@tightboxstyle={top=-6pt,bottom=-6pt}{##2}{fancy@box@example}{example},listing only,listing options={language=#1,##1}}
}

% listing style inherited by all languages
\lstset{
	numbers=none, % where to put the line-numbers
	showspaces=false, % show spaces with underscores
	showstringspaces=false, % underline spaces within strings
	showtabs=false, % show tabs within strings with underscores
	tabsize=2, % sets default tabsize to 2 spaces
	breaklines=true, % sets automatic line breaking
	breakatwhitespace=false,
	columns=fullflexible,
	commentstyle=\color{green},
	keywordstyle=\color{blue},
	stringstyle=\color{red},
}

% easy font size selection
\lstdefinestyle{small}{basicstyle=\sf\small\selectfont}
\lstdefinestyle{footnotesize}{basicstyle=\sf\footnotesize\selectfont}
\lstdefinestyle{scriptsize}{basicstyle=\sf\scriptsize\selectfont}
\lstdefinestyle{tiny}{basicstyle=\sf\tiny\selectfont}

% language-specific styles
\definelanguage{code}{
	style=small,
	language=Java,
	keywordstyle=\bf,
	moredelim=**[is][\color{red}]{@}{@},
	moredelim=**[is][\color{green}]{?}{?},
	moredelim=**[is][\color{blue}]{~}{~},
	moredelim=[is][\sout]{|}{|},
}
 % TODO merge into template?
\input{../se1/template/slidetemplate} % TODO merge into template?
% BOOKS

\newcommand{\cleancode}{\href{https://learning.oreilly.com/library/view/clean-code-a/9780136083238/}{Uncle Bob's Clean Code}}

\newcommand{\codecomplete}{\href{https://www.pearson.de/code-complete-9780735635074}{Code Complete}}

\newcommand{\cpp}{\href{https://dl.acm.org/doi/book/10.5555/518791}{C++}}

\newcommand{\dragonbook}{\href{https://dl.acm.org/doi/10.5555/1095594}{Dragon Book}}

\newcommand{\elementsofprogrammingstyle}{\href{http://www2.ing.unipi.it/~a009435/issw/extra/kp_elems_of_pgmng_sty.pdf}{Kernighan and Plauger}}

\newcommand{\ghezzi}{\href{https://www.researchgate.net/publication/220689590_Fundamentals_of_software_engineering_2_ed}{Ghezzi/Jazayeri/Mandrioli}}

\newcommand{\gof}{\gofen{} \deutsch{\gofde}}
\newcommand{\gofen}{\href{https://learning.oreilly.com/library/view/design-patterns-elements/0201633612/}{Gang of Four}}
\newcommand{\gofde}{\href{https://learning.oreilly.com/library/view/design-patterns-/9783826697005/}{GoF}}

\newcommand{\goodliffe}{\href{https://www.oreilly.com/library/view/becoming-a-better/9781491905562/}{Goodliffe 2014}}

\newcommand{\handbuch}{\href{https://www.rheinwerk-verlag.de/handbuch-fuer-softwareentwickler/}{Krypczyk/Bochkor}}

\newcommand{\javascript}{\href{https://www.oreilly.com/library/view/javascript-the-good/9780596517748/}{Crockford 2008}}

\newcommand{\lientzswanson}{\href{https://dl.acm.org/doi/abs/10.5555/601062}{Lientz and Swanson 1980}}

\newcommand{\ludewiglichter}{\href{https://dpunkt.de/produkt/software-engineering-2/}{Ludewig and Lichter}}

\newcommand{\refactoring}{\href{https://www.pearson.de/refactoring-improving-the-design-of-existing-code-9780134757599}{Fowler's Refactoring}}

\newcommand{\sommervillelink}[1]{\href{https://www.pearson.de/software-engineering-9783868943443}{#1}} % keep consistent with books/sommerville-softwareengineering.txt
\newcommand{\sommerville}{\sommervillelink{Sommerville}}

\newcommand{\thepragmaticprogrammer}{\href{https://pragprog.com/titles/tpp20/the-pragmatic-programmer-20th-anniversary-edition/}{The Pragmatic Programmer}}

% PAPERS

\newcommand{\jonesbestpractice}{\href{https://dl.acm.org/doi/10.5555/1594755}{Jones}}

\newcommand{\lehmanslaws}{\href{https://ieeexplore.ieee.org/iel3/5031/13795/00637156.pdf}{Lehman et al.\ 1997}}

\newcommand{\thehumbleprogrammer}{\href{https://dl.acm.org/doi/10.1145/1283920.1283927}{The Humble Programmer}}

% STANDARDS

\newcommand{\ieeesixten}{\href{https://ieeexplore.ieee.org/document/159342}{IEEE Std 610.12}}

\newcommand{\umlrefman}{\href{https://idsi.md/files/file/referinte_utile_studenti/The\%20Unified\%20Modeling\%20Language\%20Reference\%20Manual.pdf}{UML Reference Manual}}

\newcommand{\umlspec}{\href{https://www.omg.org/spec/UML/2.5.1/PDF}{UML 2.5.1}}

\newcommand{\umluserguide}{\href{http://patologia.com.mx/informatica/uug.pdf}{UML User Guide}}

% WEB

\newcommand{\isoiectfzoz}{\href{https://iso25000.com/index.php/en/iso-25000-standards/iso-25010}{ISO/IEC 25010}}

\newcommand{\javastyleguide}{\href{https://google.github.io/styleguide/javaguide.html}{Google 's Java Style Guide}}

\newcommand{\seeconomics}{\href{https://citeseerx.ist.psu.edu/document?repid=rep1&type=pdf&doi=970fc1889f7097bbb5c0a13c965567070a0df6d2}{SE Economics}}

% OLD LITERATURE

%\newcommand{\fospl}{\href{http://link.springer.com/book/10.1007/978-3-642-37521-7}{Apel~et~al.~2013}}
%\newcommand{\featureide}{\href{https://link.springer.com/book/10.1007\%2F978-3-319-61443-4}{Meinicke~et~al.~2013}}
%\newcommand{\icpl}{\href{https://dl.acm.org/doi/10.1145/2362536.2362547}{Johansen et al.\ 2012}}
%\newcommand{\jacobson}{\href{https://katalog.ub.uni-paderborn.de/local/r/990005758980106463?sr[q,any]=Ivar\%20Jacobson}{Jacobson 2007}} % not yet used

% literature from SPL
%\newcommand{\lehmanslaws}{\href{https://ieeexplore.ieee.org/iel3/5031/13795/00637156.pdf}{Lehman~et~al.~1997}}
%\newcommand{\ludewiglichter}{\href{https://learning.oreilly.com/library/view/-/9781457184932/?ar}{Ludewig~and~Lichter~2013}}
%\newcommand{\oosoftwareconstruction}{\href{https://dl.acm.org/doi/10.5555/261119}{Meyer~1997}}
%\newcommand{\thehumbleprogrammer}{\href{https://dl.acm.org/doi/10.1145/1283920.1283927}{The~Humble~Programmer}}

% introduced: introduction
% reused: summary
\newcommand{\slideSEvsProgramming}{
	\begin{fancycolumns}[animation=none]
		\centering
		\pic[height=60mm]{misc/ulm-muenster}
	\nextcolumn
		\centering\pic[height=60mm]{misc/tarp-tent-cropped}
	\end{fancycolumns}
	\begin{note}{}
		\centering what is needed besides programming will be motivated and shown throughout this course
	\end{note}
}

% introduced: implementation
% reused: summary, se2-compilation
\newcommand{\slideProgrammingLanguagesToday}{
	\begin{fancycolumns}
		\begin{example}{Today\mysource{\jonesbestpractice\ + \handbuch}}
			\begin{itemize}
				\item 2002: C\# by Microsoft
				\item 2009: Go by Google
				\item 2010: Rust by Mozilla Research
				\item 2014: Swift by Apple
				\item thousands of programming languages
				\item very few programming languages used for more than 10 years
				\item languages used for more than 25 years: Ada, C, C++, COBOL, Java, Objective C, PL/I, SQL, Visual Basic, \ldots
			\end{itemize}
		\end{example}
		\nextcolumn
		\begin{note}{Many Languages\mysource{\jonesbestpractice}}
			\begin{itemize}
				\item good: fit for every use case
				\item bad: developer training for new and dead languages, expensive tool support
			\end{itemize}
		\end{note}
	\end{fancycolumns}
}

% introduced: implementation
% reused: se2-compilation
\newcommand{\slideTiobeDiagram}{
	\only<1,6->{\picDark[width=\linewidth]{history/tiobe}}%
	\only<2|handout:0>{\picDark[width=\linewidth]{history/tiobe-python}}%
	\only<3|handout:0>{\picDark[width=\linewidth]{history/tiobe-cpp}}%
	\only<4|handout:0>{\picDark[width=\linewidth]{history/tiobe-c}}%
	\only<5|handout:0>{\picDark[width=\linewidth]{history/tiobe-java}}%
}
\newcommand{\slideTiobeTable}{
	\centering\picDark[width=.9\linewidth]{history/tiobe-longterm}
}

% introduced: implementation
% reused: designpatterns
\newcommand{\slideWindowsCalculator}{
	\begin{frame}{Calc on Windows 10\ \mytitlesource{\href{https://youtu.be/LKDbQfzzGJo?t=1544}{youtube.de}}}
		\begin{fancycolumns}
			\centering\picDark[width=.66\linewidth]{failures/win10-calc-scientific}
			\nextcolumn
			\centering\picDark[width=.66\linewidth]{failures/win10-calc-standard}
		\end{fancycolumns}
	\end{frame}
}

% introduced: testing
% reused: summary
\newcommand{\slideMindmapQualityAssurance}[7]{
	\vspace{-12mm}\hfill
	\begin{tikzpicture}
		\path[small mindmap,
		every node/.style={concept,font=\scriptsize},
		emph/.style={font=\bfseries\scriptsize},
		concept color=foreground!20!background,
		level 1/.append style={level distance=25mm,sibling angle=360/6},
		level 2/.append style={level distance=20mm,sibling angle=360/6},
		level 3/.append style={level distance=20mm,sibling angle=360/8},
		]
		node {Quality Assurance \deutsch{Qualitätssicherung}}
		[counterclockwise from=210]
		child[#1] { node {constructive} 
			[clockwise from=225]
			child[concept color=blue!20!background,#4] { node {Coding Guidelines} }
		}
		child[#2] { node {analytical} 
			[counterclockwise from=240]
			child[concept color=green!20!background,#5] { node {analysis}
				[counterclockwise from=180]
				child { node {Compilation} }
				child { node {Code Reviews} }
			}
			child[concept color=red!20!background,#6] { node {execution}
				[counterclockwise from=315]
				child { node {White-Box Testing} }
				child { node {Black-Box Testing} }
			}
		}
		child[#3] { node {organizational} 
			[clockwise from=-45]
			child[concept color=orange!20!background,#7] { node {Software Project Management} }
		}
		;
	\end{tikzpicture}
}

% Modification of Mindmap above (\slideMindmapQualityAssurance) for SE2
\newcommand{\slideMindmapQualityAssuranceMod}[7]{
	\vspace{-12mm}\hfill
	\begin{tikzpicture}
		\path[small mindmap,
		every node/.style={concept,font=\scriptsize},
		emph/.style={font=\bfseries\scriptsize},
		concept color=foreground!20!background,
		level 1/.append style={level distance=25mm,sibling angle=360/6},
		level 2/.append style={level distance=20mm,sibling angle=360/6},
		level 3/.append style={level distance=20mm,sibling angle=360/8},
		]
		node {Quality Assurance \deutsch{Qualitätssicherung}}
		[counterclockwise from=210]
		child[#1] { node {constructive} 
			[clockwise from=225]
			child[concept color=blue!20!background,#4] { node {Coding Guidelines} }
		}
		child[#2] { node {analytical} 
			[counterclockwise from=240]
			child[concept color=green!20!background,#5] { node {analysis}
				[counterclockwise from=180]
				child { node {Compilation} }
				child { node {Code Reviews} }
				child { node {Static Analysis} } %added for SE2
			}
			child[concept color=red!20!background,#6] { node {execution}
				[counterclockwise from=315]
				child { node {White-Box Testing} }
				child { node {Black-Box Testing} }
			}
		}
		child[#3] { node {organizational} 
			[clockwise from=-45]
			child[concept color=orange!20!background,#7] { node {Software Project Management} }
		}
		;
	\end{tikzpicture}
}

% introduced: changes
% reused: summary
\newcommand{\slideEvolutionAndMaintenance}{
	\begin{fancycolumns}[t]
		\begin{note}{Evolution}
			\begin{itemize}
				\item new or removed functionality
				\item typically larger changes
				\item often foreseen changes
				\item results in upgrades, service packs, or cumulative updates
			\end{itemize}
		\end{note}
		\begin{example}{Minor Release}
			new minor version: 2.3.1 $\Rightarrow$ 2.4.0
		\end{example}
		\begin{example}{Major Release}
			new major version: 2.3.1 $\Rightarrow$ 3.0.0
		\end{example}
		\nextcolumn
		\begin{note}{Maintenance}
			\begin{itemize}
				\item mostly corrections
				\item typically smaller changes
				\item often unforeseen changes
				\item results in patches and hot fixes
			\end{itemize}
		\end{note}
		\begin{example}{Patch Release}
			new patch version: 2.3.1 $\Rightarrow$ 2.3.2
		\end{example}
		\begin{example}{Linux Kernel Releases\mysource{\href{https://en.wikipedia.org/wiki/Linux_kernel_version_history}{wikipedia.org}}}
			\begin{itemize}
				%\item 6.13.0 -- 6.13.11
				%\item 6.12.0 -- 6.12.23
				%\item 6.11.0 (Sep 24) -- 6.11.11 (Dec 24)
				%\item 6.10.0 (Jul 24) -- 6.10.14 (Oct 24)
				\item 6.9.0 (May 24) -- 6.9.10 (Jul 24)
				\item 6.8.0 (Mar 24) -- 6.8.12 (May 24)
				\item 6.7.0 (Jan 24) -- 6.7.12 (Apr 24)
				\item 6.6.0 (Oct 23) -- 6.6.87+ (Dec 26) LTS
			\end{itemize}
		\end{example}
	\end{fancycolumns}
}

% introduced: versioncontrol
% reused: ?
\newcommand{\slideBranchingAndMerging}{
	\begin{fancycolumns}[animation=none,b,widths={58}]
		\only<1|handout:0>{\trunkbranch{}{}}%
		\only<2|handout:0>{\trunkbranch{,draw=green}{}}%
		\only<3-|handout:1>{\trunkbranch{,draw=green}{,draw=orange}}%
		\only<1->{
			\begin{note}{}
				\begin{itemize}
					\item simultaneous, independent development
					\item option to merge in the future
					\item main development on branch \texttt{main} (formerly \texttt{master})
					\item parallel developments on branches for new features, bug fixes, multiple versions, \ldots
					\item avoiding name \texttt{master} for branches:\\\texttt{git config --global init.defaultBranch main}
				\end{itemize}
			\end{note}
		}
		\nextcolumn
		\uncover<4-|handout:1>{\picDark[width=\linewidth]{versioncontrol/master-main-github}}%
	\end{fancycolumns}
}

% introduced: management
% reused: summary
\newcommand{\forwardpass}[1]{#1}
\newcommand{\backwardpass}[1]{#1}
\newcommand{\buffer}[1]{#1}
\newcommand{\networknode}[7]{
	\begin{tikzpicture}[every node/.style={draw=foreground,fill=background,anchor=base,minimum height=5mm,text width=7.5mm,align=center},inner xsep=0mm,line width=.5pt,node distance=-.5pt]
		\node[text width=22.5mm] (task) {#1};
		\node (d) [above=of task] {#3};
		\node (es) [left=of d] {\forwardpass{#2}};
		\node (ef) [right=of d] {\forwardpass{#4}};
		\node (b) [below=of task] {\buffer{#6}};
		\node (ls) [left=of b] {\backwardpass{#5}};
		\node (lf) [right=of b] {\backwardpass{#7}};
	\end{tikzpicture}
}
\newcommand{\slideGanttAndNetwork}{
	\begin{fancycolumns}
		\begin{note}{Gantt Chart}
			\begin{itemize}
				\item very common technique
				\item many tools available
				\item great visualization of timing and progress
			\end{itemize}
		\end{note}
		\begin{exampletight}{}
			\centering\picDark[width=.85\linewidth]{management/gantt-chart}
		\end{exampletight}
		\nextcolumn
		\begin{note}{Network Diagram \deutsch{Netzplan}}
			\begin{itemize}
				\item clear visualization of dependencies
				\item explicitly includes buffer times\\(cf.\ metra potential method)
			\end{itemize}
		\end{note}
		\begin{exampletight}{}
			\centering
			\begin{tikzpicture}[xscale=3,yscale=-2.1,inner sep=0,edge/.style={->,>={Stealth[round]},semithick}]
				\node[fill=red!30!background,font=\tiny] (legend) at (1,0) {\networknode{\normalsize Task}{earliest start}{duration}{earliest finish}{latest start}{buffer}{latest finish}};
				\node (intro) at (1,1) {\networknode{Introduction}{0}{1}{1}{10}{10}{11}};
				\node (background) at (0,0) {\networknode{Background}{0}{3}{3}{0}{0}{3}};
				\node (concept) at (0,1) {\networknode{Concept}{3}{4}{7}{3}{0}{7}};
				\draw[edge] (background) to (concept);
			\end{tikzpicture}
		\end{exampletight}
	\end{fancycolumns}
}

% introduced: management
% reused: process
\newcommand{\definitionuserstory}[1]{
	\begin{definition}{#1User Story \mysource{\sommerville}}
		\begin{itemize}
			\item a scenario of use that might be experienced by a system user
			\item aka.\ \emph{story card} as user stories are sometimes written on physical cards
			\item user stories are typically prioritized by the customer
			\item subset of all user stories is chosen for the next release
		\end{itemize}
	\end{definition}
}

% introduced: process
% reused: 
\newcommand{\slideWaterfallModel}{
	\begin{fancycolumns}
		\begin{definition}{Waterfall Model}
			\begin{itemize}
				\item first process model, motivated by practice
				\item by \href{https://scholar.google.de/scholar?cluster=8624196209257442707}{Winston W. Royce 1970}
				\item each development phase ends by the approval of one or more documents (document-driven process model)
				\item phases do not overlap
				\item numerous variants with varying number of phases: 5--7
				\item here: simplified variant by Sommerville
			\end{itemize}
		\end{definition}
		\nextcolumn
		\diagramWaterfallModel
	\end{fancycolumns}
}
\newcommand{\diagramWaterfallModel}{
	\begin{tikzpicture}[yscale=-1.2,xscale=.75,phase/.style={draw=foreground,thick,rounded rectangle,fill=blue!10!background,align=center,text width=21mm},label/.style={auto,bend right,align=left},to/.style={->,>={Stealth[round]},thick}]
			\node[phase,visible on=<6->] (1) at (0,0) {Requirements Analysis};
			\node[phase,visible on=<7->] (2) at (1,1) {System and Software Design};
			\node[phase,visible on=<2->] (3) at (2,2) {Implementation and Unit Testing};
			\node[phase,visible on=<4->] (4) at (3,3) {Integration and System Testing};
			\node[phase,visible on=<3->] (5) at (4,4) {Operation and Maintenance};
			\draw[to,visible on=<8->,blue] (1) -| node[label] {System\\Specification} (2.30);
			\draw[to,visible on=<8->,blue] (2) -| node[label] {Design\\Specification} (3.30);
			\draw[to,visible on=<5->,blue] (3) -| node[label] {Program\\Documentation} (4.30);
			\draw[to,visible on=<5->,blue] (4) -| node[label] {Test\\Documentation} (5.30);
			\draw[to,visible on=<9|handout:0>,red] (5) -| (1.210);
			\draw[to,visible on=<9|handout:0>,red] (5) -| (2.210);
			\draw[to,visible on=<9|handout:0>,red] (5) -| (3.210);
			\draw[to,visible on=<9|handout:0>,red] (5) -| (4.210);
			\draw[to,visible on=<10->,orange] (2) -| (1.210);
			\draw[to,visible on=<10->,orange] (3) -| (2.210);
			\draw[to,visible on=<10->,orange] (4) -| (3.210);
			\draw[to,visible on=<10->,orange] (5) -| (4.210);
	\end{tikzpicture}
}

% introduced: process
% reused: maintenance
\newcommand{\diagramVModel}{
	\begin{tikzpicture}[yscale=-1.45,xscale=.45,phase/.style={draw=foreground,thick,rounded rectangle,fill=blue!10!background,align=center,text width=17mm},label/.style={midway,anchor=west},to/.style={->,>={Stealth[round]},thick}]
			\node[phase,visible on=<2->] (1) at (0,0) {Requirements Elicitation};
			\node[phase,visible on=<3->] (2) at (1,1) {System Modeling};
			\node[phase,visible on=<4->] (3) at (2,2) {Architecture Design};
			\node[phase,visible on=<5->] (4) at (3,3) {Software Design};
			\node[phase,visible on=<6->] (5) at (6,4) {Implemen-tation};
			\node[phase,visible on=<7->] (6) at (9,3) {Unit\\Testing};
			\node[phase,visible on=<8->] (7) at (10,2) {Integration Testing};
			\node[phase,visible on=<9->] (8) at (11,1) {System Testing};
			\node[phase,visible on=<10->] (9) at (12,0) {Acceptance Testing};
			\draw[to,visible on=<3->,blue] (1) to node[label] {Requirements Specification} (2);
			\draw[to,visible on=<4->,blue] (2) to node[label] {System Specification} (3);
			\draw[to,visible on=<5->,blue] (3) to node[label] {Architecture Spec.} (4);
			\draw[to,visible on=<6->,blue] (4) to node[label] {Design Spec.} (5.165);
			\draw[to,visible on=<7->,blue] (5.15) to node[label] {Component} (6);
			\draw[to,visible on=<8->,blue] (6) to node[label] {Components} (7);
			\draw[to,visible on=<9->,blue] (7) to node[label] {System} (8);
			\draw[to,visible on=<10->,blue] (8) to node[label] {System} (9);
			\draw[to,visible on=<11->,orange,dashed] (1) to node[auto] {test cases} (9);
			\draw[to,visible on=<12->,orange,dashed] (2) to (8);
			\draw[to,visible on=<12->,orange,dashed] (3) to (7);
			\draw[to,visible on=<12->,orange,dashed] (4) to (6);
	\end{tikzpicture}
}
\newcommand{\diagramForwardEngineering}{
	\begin{tikzpicture}[yscale=-1.45,xscale=.45,phase/.style={draw=foreground,thick,rounded rectangle,fill=blue!10!background,align=center,text width=17mm},label/.style={midway,anchor=west},to/.style={->,>={Stealth[round]},thick}]
			\node[phase,visible on={<1-2,6>}] (1) at (0,0) {Requirements Elicitation};
			\node[phase,visible on={<1-3>}] (2) at (1,1) {System Modeling};
			\node[phase,visible on={<1,3-4>}] (3) at (2,2) {Architecture Design};
			\node[phase,visible on={<1,4-5>}] (4) at (3,3) {Software Design};
			\node[phase,visible on={<1,5-6>}] (5) at (4,4) {Implemen-tation};
			\draw[to,visible on={<1-2>},blue] (1) to node[label] {Requirements Specification} (2);
			\draw[to,visible on={<1,3>},blue] (2) to node[label] {System Specification} (3);
			\draw[to,visible on={<1,4>},blue] (3) to node[label] {Architecture Specification} (4);
			\draw[to,visible on={<1,5>},blue] (4) to node[label] {Design Specification} (5);
			\draw[to,visible on={<6|handout:0>},blue] (1) to node[label] {Minor Change Request} (5);
	\end{tikzpicture}
}
\newcommand{\diagramReverseEngineering}{
	\begin{tikzpicture}[yscale=-1.45,xscale=.8,phase/.style={draw=foreground,thick,rounded rectangle,fill=blue!10!background,align=center,text width=17mm},label/.style={midway,anchor=west},to/.style={->,>={Stealth[round]},thick}]
			\node[phase,visible on={<1,2,6>}] (1) at (0,0) {Requirements Elicitation};
			\node[phase,visible on={<1,4,6>}] (2) at (1,1) {System Modeling};
			\node[phase,visible on={<1,5->}] (3) at (2,2) {Architecture Design};
			\node[phase,visible on={<1,3,6>}] (4) at (3,3) {Software Design};
			\node[phase,visible on={<1->}] (5) at (4,4) {Implemen-tation};
			\draw[to,visible on={<1,6>},orange] (2) to (1);
			\draw[to,visible on={<1,6>},orange] (3) to (2);
			\draw[to,visible on={<1,6>},orange] (4) to (3);
			\draw[to,visible on={<2|handout:0>},orange] (5) to (1);
			\draw[to,visible on={<4|handout:0>},orange] (5) to (2);
			\draw[to,visible on={<5|handout:0>},orange] (5) to (3);
			\draw[to,visible on={<1,3,6>},orange] (5) to (4);
	\end{tikzpicture}
}

% introduced: process
% reused: process
\newcommand{\diagramScrum}{
	\pic[width=\linewidth]{process/scrum}
}

% introduced: requirements
% reused: summary
\newcommand{\slideMindmapNonFunctionalRequirements}{
	{\newcommand{\requirements}{requirements}%
		\begin{tikzpicture}
			\path[%grow cyclic,
			small mindmap,
			%every node/.style=concept,
			concept color=green!20!background,
			%	text width=20mm,align=flush center,
			%	level 1/.append style={level distance=27mm,sibling angle=360/3},
			]
			node[concept] {non-functional \requirements}
			[counterclockwise from=195]
			child[concept color=blue!20!background] { node[concept] {product \requirements} 
				[counterclockwise from=75]
				child[visible on=<2->] { node[concept] {usability \requirements} }
				child[visible on=<3->] { node[concept] {efficiency \requirements} 
					[counterclockwise from=165]
					child { node[concept] {performance \requirements} }
					child { node[concept] {space \requirements} }
				}
				child[visible on=<4->] { node[concept] {dependability \requirements} }
				child[visible on=<5->] { node[concept] {security \requirements} }
			}
			child[concept color=red!20!background] { node[concept] {organizational \requirements} 
				[counterclockwise from=210]
				child[visible on=<6->] { node[concept] {environmental \requirements} }
				child[visible on=<7->] { node[concept] {operational \requirements} }
				child[visible on=<8->] { node[concept] {development \requirements} }
			}
			child[concept color=orange!20!background] { node[concept] {external \requirements} 
				[counterclockwise from=285]
				child[visible on=<9->] { node[concept] {regulatory \requirements} }
				child[visible on=<10->] { node[concept] {ethical \requirements} }
				child[visible on=<11->] { node[concept] {legislative \requirements}
					[counterclockwise from=-15]
					child { node[concept] {accounting \requirements} }
					child { node[concept] {safety/security \requirements} }
				}
			}
			;
	\end{tikzpicture}}
}

% introduced: modeling
% reused: architecture, design, summary
% 
% legend: gray (will not be taught), bold (taught in this lecture), normal font (will be taught)
% 02-requirements: use case diagrams
% 03-modeling: UML, activity diagrams, state machine diagrams
% 04-architecture: component diagrams
% 05-design: class diagrams, sequence diagrams
%
% parameters: use case, activity, state machine, component, class, sequence, structure, behavior, interactions
\tikzset{notTaughtUMLDiagrams/.style={text=gray}}
\newcommand{\slideMindmapUMLdiagrams}[9]{
	{\centering
		\begin{tikzpicture}
			\path[small mindmap,
			every node/.style={concept,font=\footnotesize},
			emph/.style={font=\bfseries\footnotesize},
			blue/.style={emph,text=blue},
			red/.style={emph,text=red},
			concept color=green!20!background,
			level 1/.append style={level distance=35mm,sibling angle=360/3},
			level 2/.append style={level distance=20mm,sibling angle=360/8},
			level 3/.append style={level distance=20mm,sibling angle=360/8},
			]
			node {UML Diagrams}
			[clockwise from=150]
			child[concept color=blue!20!background,#7] { node {Structure Diagrams} 
				[counterclockwise from=15]
				child { node[#4] {Component Diagram} }
				child { node[#5] {Class Diagram} }
				child[notTaughtUMLDiagrams] { node {Profile Diagram} }
				child[notTaughtUMLDiagrams] { node {Composite Structure Diagram} }
				child[notTaughtUMLDiagrams] { node {Deployment Diagram} }
				child[notTaughtUMLDiagrams] { node {Object Diagram} }
				child[notTaughtUMLDiagrams] { node {Package Diagram} }
			}
			child[concept color=red!20!background,#8] { node {Behavior Diagrams} 
				[clockwise from=135]
				child { node[#1] {Use Case Diagram} }
				child { node[#2] {Activity Diagram} }
				child { node[#3] {State Machine Diagram} }
				child[concept color=orange!20!background,#9] { node {Interaction Diagrams} 
					[clockwise from=30]
					child { node[#6] {Sequence Diagram} }
					child[notTaughtUMLDiagrams] { node {Communi-cation Diagram} }
					child[notTaughtUMLDiagrams] { node {Interaction Overview Diagram} }
					child[notTaughtUMLDiagrams] { node {Timing Diagram} }
				}
			}
			;
	\end{tikzpicture}}
}

% introduced: architecture
% reused: summary
\newcommand{\slideArchitecturalPattern}{
	\begin{fancycolumns}
		\centering
		\begin{definition}{Architectural Pattern\mysource{\sommerville}}
			\mycite{Architectural patterns capture the essence of an architecture that has been used in different software systems. [...] Architectural patterns are a means of reusing knowledge about generic system architectures.}
		\end{definition}
		\nextcolumn
		\begin{note}{Goals}
			\partofpage{35}{
				\pic[width=\textwidth]{architecture/database}
			}
			\partofpage{60}{
				\begin{itemize}
					\item preserve knowledge of software architects
					\item reuse of established architectures
					\item enable efficient communication
				\end{itemize}
			}
		\end{note}
	\end{fancycolumns}
}

% introduced: architecture
% reused: 
\newcommand{\slidePipeAndFilter}{
	\begin{fancycolumns}[animation=none]
		\begin{definition}{Pipe-and-Filter Architecture \mysource{\sommerville}}
			\begin{itemize}
				\item \emph{Problem}: data is processed in numerous processing steps, which are prone to change
				\item \emph{Idea}: modularization of each processing step into a component
				\item filter components process a stream of data continously
				\item pipes transfer data unchanged from filter output to filter input
			\end{itemize}
		\end{definition}
		\uncover<2->{
			\begin{example}{Pipe Operator in UNIX}
				\mycite{\texttt{ls -al | grep '2020' | grep -v 'Nov' | more}} searches files in a folder from the year 2020 except those from November and delivers the results in pages.
			\end{example}
		}
		\nextcolumn%
		\uncover<3->{\hfill\picDark[width=.8\textwidth]{architecture/pipe-and-filter}}%
	\end{fancycolumns}
}

% introduced: se1-reuse
% reused: se2-designpatterns
\newcommand{\slideGangOfFour}{
	\begin{frame}{\insertsubsection\ \mytitlesource{\gof}}
	\begin{fancycolumns}[animation=none]
		\centering\pic[width=.66\linewidth]{books/gof}
		\nextcolumn
		\centering\pic[height=27mm,trim=185 0 184 0,clip]{people/erich-gamma}
		\pic[height=27mm]{people/richard-helm}
		
		\pic[height=27mm,trim=0 0 5 0,clip]{people/ralph-johnson}
		\pic[height=27mm,trim=0 0 160 0,clip]{people/john-vlissides}
	\end{fancycolumns}
\end{frame}
}

% introduced: se1-reuse
% reused: se2-designpatterns
\newcommand{\slideDesignPatterns}{
	\begin{frame}{\insertsubsection\ \mytitlesource{\gofen}}
		\begin{fancycolumns}
			\begin{note}{Motivation}
				\mycite{Designing object-oriented software is hard, and designing \emph{reusable} object-oriented software is even harder. [...]  It takes a long time for novices to learn what good object-oriented design is all about. Experienced designers evidently know something inexperienced ones don't. What is it?}
			\end{note}
			\nextcolumn
			\begin{definition}{Design Patterns \deutsch{Entwurfsmuster}}
				\setlength\tabcolsep{1mm}
				\begin{tabularx}{\textwidth}{rX}				
					pattern name & for communication and high-level abstraction\\
					problem & when to apply the pattern\\
					solution & template on how to arrange classes and objects\\
					consequences & trade-offs of applying the pattern
				\end{tabularx}
			\end{definition}
			\begin{note}{Kinds of Patterns}
				\mycite{Creational patterns \deutsch{Erzeugungsmuster} concern the process of object creation. Structural patterns \deutsch{Strukturmuster} deal with the composition of classes or objects. Behavioral patterns \deutsch{Verhaltensmuster} characterize the ways in which classes or objects interact and distribute responsibility.}
			\end{note}
		\end{fancycolumns}
	\end{frame}
}

% introduced: se1-reuse
% reused: se2-designpatterns
%
% legend: gray (will not be taught), bold (taught in this lecture), normal font (will be taught)
% 
% parameters: object adapter, composite, decorator, singleton, abstract factory, observer, structural, creational, behavioral
\tikzset{notTaughtDesignPatterns/.style={text=gray}}
\newcommand{\slideMindmapDesignPatterns}[9]{
	{\centering
		\begin{tikzpicture}
			\path[small mindmap,
			every node/.style={concept,font=\scriptsize},
			emph/.style={font=\bfseries\scriptsize},
			blue/.style={emph,text=blue},
			red/.style={emph,text=red},
			concept color=green!20!background,
			level 1/.append style={level distance=35mm,sibling angle=65},
			level 2/.append style={level distance=20mm,sibling angle=360/12},
			level 3/.append style={level distance=20mm,sibling angle=360/8},
			]
			node {Design Patterns}
			[clockwise from=155]
			child[concept color=blue!20!background,#7] { node {Structural Patterns} 
				[clockwise from=215]
				child[#1] { node {Object Adapter} }
				child[notTaughtDesignPatterns] { node {Class Adapter} }
				child[notTaughtDesignPatterns] { node {Bridge} }
				child[#2] { node {Composite} }
				child[#3] { node {Decorator} }
				child[notTaughtDesignPatterns] { node {Facade} }
				child[notTaughtDesignPatterns] { node {Flyweight} }
				child[notTaughtDesignPatterns] { node {Proxy} }
			}
			child[concept color=red!20!background,#8] { node {Creational Patterns} 
				[clockwise from=150]
				child[#5] { node {Abstract Factory} }
				child[notTaughtDesignPatterns] { node {Builder} }
				child[notTaughtDesignPatterns] { node {Factory Method} }
				child[notTaughtDesignPatterns] { node {Prototype} }
				child[#4] { node {Singleton} }
			}
			child[concept color=orange!20!background,#9] { node {Behavioral Patterns} 
				[clockwise from=175]
				child[notTaughtDesignPatterns] { node {Chain of Responsibility} }
				child[notTaughtDesignPatterns] { node {Command} }
				child[notTaughtDesignPatterns] { node {Interpreter} }
				child[notTaughtDesignPatterns] { node {Iterator} }
				child[notTaughtDesignPatterns] { node {Mediator} }
				child[notTaughtDesignPatterns] { node {Memento} }
				child[#6] { node {Observer} }
				child[notTaughtDesignPatterns] { node {State} }
				child[notTaughtDesignPatterns] { node {Strategy} }
				child[notTaughtDesignPatterns] { node {Template Method} }
				child[#6] { node {Visitor} } % TODO check how to parameterize that for SE1
			}
			;
	\end{tikzpicture}}
}

% design patterns as uml class diagrams modeled with diagrams.net
\newcommand{\objectadapter}[1]{\picWhite[#1,page=1,trim=110 240 495 200,clip]{design/diagramsnet-designpatterns}}
\newcommand{\composite}[1]{\picWhite[#1,page=2,trim=200 300 405 115,clip]{design/diagramsnet-designpatterns}}
\newcommand{\compositeexample}[1]{\picWhite[#1,page=3,trim=200 100 400 100,clip]{design/diagramsnet-designpatterns}}
\newcommand{\decorator}[1]{\picWhite[#1,page=4,trim=185 205 355 135,clip]{design/diagramsnet-designpatterns}}
\newcommand{\decoratorexample}[1]{\picWhite[#1,page=5,trim=185 260 340 135,clip]{design/diagramsnet-designpatterns}}

\newcommand{\singleton}[1]{\picWhite[#1,page=6,trim=55 425 665 85,clip]{design/diagramsnet-designpatterns}}
\newcommand{\singletonexample}[1]{\picWhite[#1,page=7,trim=55 410 665 85,clip]{design/diagramsnet-designpatterns}}
\newcommand{\abstractfactory}[1]{\picWhite[#1,page=8,trim=200 355 280 65,clip]{design/diagramsnet-designpatterns}}

\newcommand{\observer}[1]{\picWhite[#1,page=9,trim=285 350 320 55,clip]{design/diagramsnet-designpatterns}}
\newcommand{\visitor}[1]{\picWhite[#1,page=10,trim=80 205 420 85,clip]{design/diagramsnet-designpatterns-new}}

% TODO integrate separate file for the calculator example
\newcommand{\profcalculator}[1]{\picWhite[#1,trim=190 185 130 105,clip]{design/diagramsnet-designpatterns}}
\newcommand{\profcalculatorcollaboration}[1]{\picWhite[#1,trim=70 0 20 90,clip]{design/diagramsnet-designpatterns-collaborations}}

% introduced: se2-02-evolution
% reused: se2-05-compilation
\newcommand{\slideSimplicityOverPerformance}{
	\begin{fancycolumns}
		\pic[width=\linewidth,trim=0 50 0 30,clip]{people/wes-dyer}
		\vspace{-7mm}
		
		\begin{note}{Wes Dyer \mysource{\href{https://twitter.com/CodeWisdom/status/801456038008520705}{twitter.com}}}
			\mycite{Make it correct, make it clear, make it concise, make it fast. In that order.}
		\end{note}
		% professor? lecturer?
		\nextcolumn
		\pic[width=\linewidth,trim=100 0 0 25,clip]{people/joshua-bloch}
		\vspace{-7mm}
		
		\begin{note}{Joshua J.\ Bloch (born 1961) \mysource{\href{https://twitter.com/codewisdom/status/876878006278598656}{twitter.com}}}
			\mycite{The cleaner and nicer the program, the faster it's going to run. And if it doesn't, it'll be easy to make it fast.}
		\end{note}
		% known for design of Java platform
	\end{fancycolumns}
}

\newcommand{\figcompilerarchfull}{
	\begin{tikzpicture}[every node/.style = {align = center, text depth = 0pt}]       
		\node[fill = background] (A) {Source Code\\as String};
		
		\node[below = 20pt of A, minimum width=2.25cm, draw, fill=background] (An1) {Scanner};
		\node[below = of An1, minimum width=2.25cm, draw, fill =background] (An2) {Parser};
		\node[below = of An2, minimum width=2.25cm, draw, fill = background] (An3) {Name and Type\\Analysis};        
		
		\begin{scope}[on background layer]
			\fill[fill = blue!25!background] ([xshift = -5pt, yshift = 5pt]An1.north west) rectangle node (H1) {} ([xshift = 5pt, yshift = -5pt]An3.south east);
		\end{scope}
		
		\node [left = 1.4cm of H1.center, color = blue, rotate=90, anchor = south, fill=background, inner sep = 0pt] {Analysis};
		
		\path[-latex]
		(A) edge ([yshift = 5pt]An1.north)
		(An1) edge node[align = right, left, font = \small] {Token\\Stream} (An2)           
		(An2) edge node[align = right, left, font = \small] {Syntax\\Tree} (An3);
		
		\node[right = 1.5cm of A, fill = background] (tgt) {Target Code\\as String};
		
		
		\node[below = 20pt of tgt, minimum width=2.25cm, draw, fill=background] (syn2) {Code Generator};
		\node[below = of syn2, minimum width=2.25cm, draw, fill=background] (syn1) {Translator};
		
		\node[below = of syn1, fill = background] (ast) {Abstract\\SyntaxTree};
		
		\path[]
		([yshift = -5pt]An3.south) edge ([yshift = -10pt] An3.south)
		([yshift = -10pt] An3.south) edge[-latex, to path={-| (\tikztotarget)}] (ast);  
		
		\path[-latex]
		(ast) edge ([yshift = -5pt] syn1.south)
		([xshift = -15pt]syn1.north) edge node[align = left, right, font = \small] (il) {Intermediate\\Language} ([xshift = -15pt]syn2.south)
		([yshift = 5pt]syn2.north) edge (tgt);
		;
		
		\begin{scope}[on background layer]
			\fill[fill = red!25!background] ([xshift = -5pt, yshift = 5pt]syn2.north west) rectangle node (H2) {} ([xshift = 5pt, yshift = -5pt]syn1.south east);
		\end{scope}
		
		\node [left = 1.4cm of H2.center, color = red, rotate=90, anchor = south, inner sep=0pt, fill=background] {Synthesis};
		
		
		\node[right = 20pt of il.east, fill = background] (op2) {Attribution \&\\Optimization};
		
		\draw (intersection cs:
		first line={(op2.center) -- ([yshift = 3cm] op2.center)},
		second line={(tgt.center) -- ([xshift = 3cm] tgt.center)}) node[fill = background] (op1) {Peep Hole\\Optimization}; 
		
		\draw (intersection cs:
		first line={(op2.center) -- ([yshift = 3cm] op2.center)},
		second line={(ast.center) -- ([xshift = 3cm] ast.center)}) node[fill = background] (op3) {Attribution \&\\Optimization}; 
		
		\path[]         
		([xshift = -18pt, yshift = 3pt]op1.west) edge[-latex] ([xshift = -3pt, yshift = 3pt]op1.west)   
		([xshift = -18pt, yshift = -3pt]op1.west) edge[latex-] ([xshift = -3pt, yshift = -3pt]op1.west) 
		
		([xshift = -18pt, yshift = 3pt]op2.west) edge[-latex] ([xshift = -3pt, yshift = 3pt]op2.west)   
		([xshift = -18pt, yshift = -3pt]op2.west) edge[latex-] ([xshift = -3pt, yshift = -3pt]op2.west) 
		
		([xshift = -18pt, yshift = 3pt]op3.west) edge[-latex] ([xshift = -3pt, yshift = 3pt]op3.west)   
		([xshift = -18pt, yshift = -3pt]op3.west) edge[latex-] ([xshift = -3pt, yshift = -3pt]op3.west)     
		;       
	\end{tikzpicture}       
}

\newcommand{\figcompilerarchnoop}{
	\begin{tikzpicture}[every node/.style = {align = center, text depth = 0pt}]       
		\node[fill = background] (A) {Source Code\\as String};
		
		\node[below = 20pt of A, minimum width=2.25cm, draw, fill=background] (An1) {Scanner};
		\node[below = of An1, minimum width=2.25cm, draw, fill = background] (An2) {Parser};
		\node[below = of An2, minimum width=2.25cm, draw, fill = background] (An3) {Name and Type\\Analysis};        
		
		\begin{scope}[on background layer]
			\fill[fill = blue!25!background] ([xshift = -5pt, yshift = 5pt]An1.north west) rectangle node (H1) {} ([xshift = 5pt, yshift = -5pt]An3.south east);
		\end{scope}
		
		\node [left = 1.4cm of H1.center, color = blue, rotate=90, anchor = south, fill=background, inner sep = 0pt] {Analysis};
		
		\path[-latex]
		(A) edge ([yshift = 5pt]An1.north)
		(An1) edge node[align = right, left, font = \small] {Token\\Stream} (An2)           
		(An2) edge node[align = right, left, font = \small] {Syntax\\Tree} (An3);
		
		\node[right = 1.5cm of A, fill = background] (tgt) {Target Code\\as String};
		
		
		\node[below = 20pt of tgt, minimum width=2.25cm, draw, fill=background] (syn2) {Code Generator};
		\node[below = of syn2, minimum width=2.25cm, draw, fill=background] (syn1) {Translator};
		
		\node[below = of syn1, fill = background] (ast) {Abstract\\SyntaxTree};
		
		\path[]
		([yshift = -5pt]An3.south) edge ([yshift = -10pt] An3.south)
		([yshift = -10pt] An3.south) edge[-latex, to path={-| (\tikztotarget)}] (ast);  
		
		\path[-latex]
		(ast) edge ([yshift = -5pt] syn1.south)
		([xshift = -15pt]syn1.north) edge node[align = left, right, font = \small] (il) {Intermediate\\Language} ([xshift = -15pt]syn2.south)
		([yshift = 5pt]syn2.north) edge (tgt);
		;
		
		\begin{scope}[on background layer]
			\fill[fill = red!25!background] ([xshift = -5pt, yshift = 5pt]syn2.north west) rectangle node (H2) {} ([xshift = 5pt, yshift = -5pt]syn1.south east);
		\end{scope}
		
		\node [left = 1.4cm of H2.center, color = red, rotate=90, anchor = south, inner sep=0pt, fill=background] {Synthesis};
	\end{tikzpicture}       
}

\newcommand{\figcompilerarchbase}{
	\begin{tikzpicture}[every node/.style = {align = center, text depth = 0pt}]       
		\node[fill = background] (A) {Source Code\\as String};
		
		\node[below = 20pt of A, minimum width=2.25cm] (An1) {\phantom{Scanner}};
		\node[below = of An1, minimum width=2.25cm] (An2) {\phantom{Parser}};
		\node[below = of An2, minimum width=2.25cm] (An3) {\phantom{Name and Type}\\\phantom{Analysis}};        
		
		\begin{scope}[on background layer]
			\fill[fill = blue!25!background] ([xshift = -5pt, yshift = 5pt]An1.north west) rectangle node (H1) {} ([xshift = 5pt, yshift = -5pt]An3.south east);
		\end{scope}
		
		\node [left = 1.4cm of H1.center, color = blue, rotate=90, anchor = south, fill=background, inner sep = 0pt] {Analysis};
		
		\path[-latex]
		(A) edge ([yshift = 5pt]An1.north);
		
		\node[right = 1.5cm of A, fill = background] (tgt) {Target Code\\as String};
		
		\node[below = 20pt of tgt, minimum width=2.25cm] (syn2) {\phantom{Code Generator}};
		\node[below = of syn2, minimum width=2.25cm] (syn1) {\phantom{Translator}};
		
		\node[below = of syn1, fill = background] (ast) {Abstract\\SyntaxTree};
		
		\path[]
		([yshift = -5pt]An3.south) edge ([yshift = -10pt] An3.south)
		([yshift = -10pt] An3.south) edge[-latex, to path={-| (\tikztotarget)}] (ast);  
		
		\path[-latex]
		(ast) edge ([yshift = -5pt] syn1.south)
		([yshift = 5pt]syn2.north) edge (tgt)
		;
		
		\begin{scope}[on background layer]
			\fill[fill = red!25!background] ([xshift = -5pt, yshift = 5pt]syn2.north west) rectangle node (H2) {} ([xshift = 5pt, yshift = -5pt]syn1.south east);
		\end{scope}
		
		\node [left = 1.4cm of H2.center, color = red, rotate=90, anchor = south, inner sep=0pt, fill=background] {Synthesis};
	\end{tikzpicture}    
}

\newcommand{\figCompiler}{
	\begin{tikzpicture}[every node/.style = {align = center, fill = background}, node distance = 1cm and 0.5cm]       
		\node[] (A) {Source Code};
		
		\node[right = of A, fill = black, text=white, minimum width=2cm] (Compiler) {Compiler};
		
		\node[right = of Compiler] (tgt) {Target Code};
		
		\path[-latex]
		(A) edge (Compiler)
		(Compiler) edge (tgt);
	\end{tikzpicture}       
}

\newcommand{\figInterpreter}{
	\begin{tikzpicture}[every node/.style = {align = center, fill = background}]     
		\node[] (A) {Source Code};
		
		\node[below = of A] (B) {Input Data};
		
		\path (A.center) -- node (H) {} (B.center);
		
		\node[right = of H, fill = black, text=white, minimum width=2cm] (Compiler) {Interpreter};
		
		\node[right = 0.5cm of Compiler] (tgt) {Output Data};
		
		\path[-latex]
		(A) edge[out = -90, in = 180] ([yshift = 2pt]Compiler.west)
		(B) edge[out = 90, in = 180] ([yshift = -2pt]Compiler.west)
		(Compiler) edge (tgt);
	\end{tikzpicture}       
}

\newcommand{\figCompilerWithAnalysisOut}{
	\begin{tikzpicture}[every node/.style = {align = center, fill = background}] 
		
		\node[fill = black, text=white, minimum width=2cm] (Compiler) {Compiler};
		
		\node[left = of Compiler] (H1) {};		 
		\node[right = of Compiler] (H2) {};
		
		\node[above = 10pt of H1.center] (A) {Source Code};
		
		\node[above = 10pt of H2.center] (X) {Target Code};
		\node[below = 10pt of H2.center] (Y) {Analysis Result};
		
		\path[-latex, looseness = 0.7]
		(A) edge[out = -90, in = 180] ([yshift = 2pt]Compiler.west)
		([yshift = 2pt]Compiler.east) edge[out = 0, in = -90] (X)
		([yshift = -2pt]Compiler.east) edge[out = 0, in = 90] (Y);
	\end{tikzpicture}       
}

\newcommand{\figInterpreterWithAnalysisOut}{
	\begin{tikzpicture}[every node/.style = {align = center, fill = background}]      
		
		\node[fill = black, text=white, minimum width=2cm] (Compiler) {Interpreter};
		
		\node[left = of Compiler] (H1) {};		 
		\node[right = of Compiler] (H2) {};
		
		\node[above = 10pt of H1.center] (A) {Source Code};
		\node[below = 10pt of H1.center] (B) {Input Data};
		
		\node[above = 10pt of H2.center] (X) {Output Data};
		\node[below = 10pt of H2.center] (Y) {Analysis Result};
		
		\path[-latex, looseness = 0.7]
		(A) edge[out = -90, in = 180] ([yshift = 2pt]Compiler.west)
		(B) edge[out = 90, in = 180] ([yshift = -2pt]Compiler.west)
		([yshift = 2pt]Compiler.east) edge[out = 0, in = -90] (X)
		([yshift = -2pt]Compiler.east) edge[out = 0, in = 90] (Y);
	\end{tikzpicture}       
}

\newcommand{\drawFileSymbol}[3][]{
	
	\fill[draw, fill=background,#1] (#2,#3) -- (#2+0.75,#3) -- (#2+1, #3-0.25) -- (#2+1, #3-1.5) -- (#2, #3-1.5) -- (#2,#3) (#2+0.75, #3) -- (#2+0.75, #3-0.25) -- (#2+1, #3-0.25) (#2+0.2, #3-0.6) -- (#2+0.8, #3-0.6) (#2+0.2, #3-0.8) -- (#2+0.8, #3-0.8) (#2+0.2, #3-1) -- (#2+0.8, #3-1);
}

\newcommand{\figJIT}{
	\begin{tikzpicture}[every node/.style = {text depth = 0pt}]
		
		\node[draw = blue, anchor = north west, fill = blue!30!background] (F) {
			\begin{tikzpicture}
				\node[minimum width=5cm, minimum height = 1cm, fill = gray!10!background, draw = gray!30!background] (HW) {};
				\node[below left = 0pt and 0pt of HW.north east, font = \small, align=right] {Real Machine /\\Hardware};
				
				\node[right = 5pt of HW.west, fill = background, draw, inner sep=4pt, rounded corners=4pt] (MC) {Machine Code};
				
				\node[above = 20pt of MC.north, align=center, fill = background, draw, inner sep=4pt] (Y) {JIT Translator};	
				
				\node[above = 14pt of Y.north, align=center, fill = background, draw, rounded corners = 4pt, inner sep=4pt] (X) {Java Byte\\Code};
				
				\path[-latex] 
				(X) edge (Y)
				(Y) edge ([yshift = 7pt]MC.north)
				;
			\end{tikzpicture}
		};
		
		\node[below left = 0pt and 0pt of F.north east, font = \small, align=right] {Java Virtual Machine};
		
		\node[above = 15pt of F.north, align=center, fill = background, draw, rounded corners = 4pt, inner sep=4pt, xshift = 1.5cm] (ID) {Input Data};
		
		
		\node[above = 13pt of F, xshift = -1.5cm] (D) {
			\begin{tikzpicture}
				\drawFileSymbol[scale = 0.5]{0.9}{0.45}
				\drawFileSymbol[scale = 0.5]{0.6}{0.3}
				\drawFileSymbol[scale = 0.5]{0.3}{0.15}
				\drawFileSymbol[scale = 0.5]{0}{0}
			\end{tikzpicture}
		};
		
		\node[yshift = -8pt, align=center, fill = background, draw, rounded corners = 4pt, inner sep=4pt] at (D) (E) {Class Files};
		
		
		\node[above = 15pt of D, align = center, draw, inner sep=4pt, fill = background] (C) {Analyzer \&\\Translator};
		\node[left = 2pt of C, rotate=90, anchor=south, fill = background] {\texttt{javac}};
		
		\node[above = 38pt of ID] (A) {
			\begin{tikzpicture}
				\drawFileSymbol[scale = 0.5]{0.9}{0.45}
				\drawFileSymbol[scale = 0.5]{0.6}{0.3}
				\drawFileSymbol[scale = 0.5]{0.3}{0.15}
				\drawFileSymbol[scale = 0.5]{0}{0}
			\end{tikzpicture}
		};
		\node[yshift = -10pt, align=center, fill = background, draw, rounded corners = 4pt, inner sep=4pt] at (A) (B) {Java Source\\Code Unit};
		
		\node[left = 15pt of F.west, align=center, fill = background, draw, rounded corners = 4pt, inner sep=4pt] (OD) {Output Data};
		
		\path[-latex]
		(B) edge (C)
		(C) edge (D)
		(E) edge ([xshift = -1.5cm]F.north)
		(ID) edge ([xshift = 1.5cm]F.north)
		(F) edge (OD)
		;
		
	\end{tikzpicture}
}

% introduced: se2-introduction
% reused: se2-dynamic
\newcommand{\slideCrowdStrike}{
	\begin{fancycolumns}[b]
		\begin{definition}{\insertsubsection\mysource{\href{https://en.wikipedia.org/wiki/CrowdStrike}{wikipedia.org}, \href{https://www.bbc.com/news/articles/c23k4yyjxp3o}{bbc.com}}}
			\begin{itemize}
				\item CrowdStrike: American cybersecurity company
				\item Falcon: Windows kernel extension to prevent against threats
				\item about 10 patches released each days
				\item update on July 19, 2024 affected 8.5 million devices running Windows 10 or 11
				\item blue screen of death
				\item no restart possible without manual intervention
				\item largest IT outage
				\item about 5,000 (5 \%) flight cancelled
				\item Delta Airlines lost 500 million dollar
				\item overall financial damage of 10 billion dollar
				\item similar problem affected Linux devices already in April 2024
			\end{itemize}
		\end{definition}
		\nextcolumn
		\small\vspace{-20mm}
		\picDark[width=\linewidth]{failures/crowdstrike}
		\vspace{-4mm}
		\begin{note}{Opening in US Congress\mysource{\href{https://www.bbc.com/news/articles/c23k4yyjxp3o}{bbc.com}}}
			\mycite{A global IT outage that impacts \emph{every sector of the economy} is a catastrophe that we would expect to see in a movie [\ldots] we would expect to be carefully executed by a malicious and sophisticated nation-state actor [but] the largest IT outage in history was due to a \emph{mistake}.}
		\end{note}
		\begin{note}{Apology in US Congress\mysource{\href{https://www.bbc.com/news/articles/c23k4yyjxp3o}{bbc.com}}}
			\mycite{On behalf of everyone at CrowdStrike, I want to apologize. We are \emph{deeply sorry} and we are determined to prevent this from ever happening again. [\ldots] I want to underscore that this was \emph{not a cyber attack}.}
		\end{note}
	\end{fancycolumns}
}

% introduced: se2-introduction
% reused: se2-designbycontract
\newcommand{\slideArianeFailure}{
	\begin{fancycolumns}[widths={35}]
		\pic[width=\linewidth]{failures/ariane5}
		\nextcolumn
		\centering\pic[width=.5\linewidth]{failures/ariane5-explosion} {\tiny\copyright{} ESA} % TODO how to add copyright? ESA
		\begin{definition}{Ariane 5 Software Failure \mysource{\lessonsofariane}}
			\begin{itemize}
				\item Ariane 5 is a European space launch vehicle
				\item 117 launches between 1996 and 2023
				\item maiden flight in June 1996
				\item explosion about 40 seconds after takeoff, no personal injury
				\item software component of Ariane 4 was reused
				\item value exceeding 16-bit signed integer due to Ariane 5 trajectory
				\item damage of 500 million dollar + one year postponed research
			\end{itemize}
		\end{definition}
	\end{fancycolumns}
}

% introduced: se2-introduction
% reused: se2-automotive
\newcommand{\slideXiaomiCrash}{
	\begin{fancycolumns}[widths={45},animation=none]
		\begin{definition}{Xiaomi SU7 Crash\mysource{\href{https://www.nytimes.com/2025/04/01/business/xiaomi-driverless-electric-car-china.html}{nytimes.com}}}
			\begin{itemize}
				\item 200,000+ sold cars in first year
				\item promoted with assisted-driving: \mycite{car can automatically accelerate or decelerate, change lanes, enter and exit a highway, and avoid construction}
				\item accident on March 29, 2025 (Saturday)
				\item car was driving at 116 km/h
				\item advanced driver assistance was turned on
				\item road closed due to reconstruction and traffic rerouted to reverse lane
				\item electric car crashed into highway guardrail with 97 km/h and caught fire
				\item three female students killed
				\item Xiaomi's stock reduced by 9\%
			\end{itemize}
		\end{definition}
		\nextcolumn
		\vspace{-10mm}
		
		\small\centering\pic[width=.7\linewidth]{automotive/xiaomi-su7}
		\begin{note}{Log of the Car\mysource{\href{https://weibo.com/5136788508/5150602603790529}{weibo.com}}}
			22:27:17 NOA activation, speed 116km/h
			
			22:28:17 Mild distraction alarm
			
			22:36:48 NOA issued a warning prompt "please hold the steering wheel"
			
			22:44:24 NOA issues a risk alert "Please note that there are obstacles ahead", make a deceleration request, and start to slow down.
			
			22:44:25 NOA was taken over and entered the human-driven state, the steering wheel turned to the left 22.0625 degrees, the brake pedal was turned on 31\%.
			
			22:44:26 Turn right steering wheel 1.0625 degrees, brake pedal open 38\%
			
			22:44:26-28 Collision between vehicle and cement guardrail
			
			22:44:28 Ecall trigger
		\end{note}
	\end{fancycolumns}
}

% introduced: se1-testing
% reused: se2-dynamic
\newcommand{\slideSoftwareQuality}{
	\begin{fancycolumns}
		\begin{definition}{Quality \mysource{\ludewiglichter}}
			Quality is the entirety of properties and characteristics of a product or process that indicate adequacy with respect to given requirements.
		\end{definition}
		\begin{definition}{Quality Assurance \mysource{\ludewiglichter}}
			Quality assurance \deutsch{Qualitätssicherung} are all activities with the goal to improve the quality.
		\end{definition}
		\nextcolumn
		\begin{note}{Expectations on Quality \mysource{\sommerville}}
			\mycite{Because of their previous experiences with buggy, unreliable software, users sometimes have low expectations of software quality. They are not surprised when their software fails. When a new system is installed, users may tolerate failures because the benefits of use outweigh the costs of failure recovery. However, as a software product becomes more established, users expect it to become more reliable.\uncover<3->{ [...] If a software product or app is very cheap, users may be willing to tolerate a lower level of reliability.}\uncover<4->{ [...] Customers may be willing to accept the software, irrespective of problems, because the costs of not using the software are greater than the costs of working around the problems.}}
		\end{note}
	\end{fancycolumns}
}


% introduced: se1-testing
% reused: se2-dynamic
\newcommand{\slideSoftwareTesting}{
\begin{fancycolumns}[height=1cm,animation=none]
	\uncover<1->{
		\begin{definition}{Software Testing \mysource{\sommerville}}
			\mycite{Testing is intended to show that a program does what it is intended to do and to discover program defects before it is put into use.}
		\end{definition}
	}
	\uncover<2->{
		\begin{definition}{Validation Testing \mysource{\sommerville}}
			\mycite{Demonstrate to the developer and the customer that the software meets its requirements.}
		\end{definition}
	}
	\uncover<3->{
		\begin{definition}{Defect Testing \mysource{\sommerville}}
			\mycite{Find inputs or input sequences where the behavior of the software is incorrect, undesirable, or does not conform to its specification.}
		\end{definition}
	}
	\nextcolumn
	%\vspace{-12mm}
	\uncover<4->{
		\begin{note}{V\&V \mysource{\seeconomics}}
			\mycite{\emph{Validation}: Are we building the right product?\\\emph{Verification}: Are we building the product right?}
		\end{note}	
	} 
	% TODO better visualization of V&V (see Inas slides). move to V model?
	\vspace{-.7mm}
	\uncover<5->{
		\begin{note}{Stages of Testing \mysource{\sommerville}}
			\begin{itemize}
				\setlength\itemsep{.1em}
				\item[1.] \mycite{\emph{Development testing}, where the system is tested during development to discover
					bugs and defects}
				\item[2.] \mycite{\emph{Release testing}, where a separate testing team tests a complete version of the
					system before it is released to users}
				\item[3.] \mycite{\emph{User testing}, where users or potential users of a system test the system in their
					own environment}
			\end{itemize}
		\end{note}
	}
	\vspace{-.7mm}
	\uncover<6->{
		\begin{note}{}
			\mycite{In \emph{manual testing}, a tester runs the program with some test data and
				compares the results to their expectations. [...] In \emph{automated testing}, the tests are encoded in a program that is run each time the system under development is to be tested.} \mysource{\sommerville}
		\end{note}
	}
\end{fancycolumns}
}

% introduced: se1-testing
% reused: se2-dynamic
\newcommand{\explTestCases}{
	\uncover<1->{
		\begin{definition}{Systematic Test \mysource{\ludewiglichter}}
			A systematic test is a test, in which
			\begin{itemize}
				\item[1.] the setup is defined,
				\item[2.] the inputs are chosen systematically,
				\item[3.] the results are documented and evaluated by criteria being defined prior to the test. 
			\end{itemize}
		\end{definition}
	}
	\uncover<2->{
		\begin{definition}{Test Case \mysource{\ludewiglichter}}
			In a test, a number of test cases are executed, whereas each test case consists \emph{input values} for a single execution and \emph{expected outputs}. An \emph{exhaustive test} refers a test in which the test cases exercise all the possible inputs.
		\end{definition}
	}
}

% introduced: se1-testing
% reused: se2-dynamic
\newcommand{\figTestDesign}{
\\\vspace{-7mm}
\hfill
\begin{tikzpicture}
	\path[small mindmap,
	every node/.style={concept,font=\scriptsize},
	emph/.style={font=\bfseries\scriptsize},
	concept color=foreground!20!background]
	node {Test-Case Design}
	child[concept color=blue!20!background,visible on={<2->}] { node {experience-based} }
	child[concept color=green!20!background,visible on={<3->}] { node {structure-based} 
		child { node {white-box testing} }
	}
	child[concept color=orange!20!background,visible on={<4->}] { node {specification-based} 
		child { node {black-box testing} }
	}
	;
\end{tikzpicture}
}

% introduced: se1-development
% reused: se2-dynamic
\newcommand{\slideStagesTesting}{
\begin{fancycolumns}[animation=none]
	\begin{definition}{1. Unit Testing \deutsch{Komponententest}}
		\uncover<2->{\begin{itemize}
				\item each component is tested independently
				\item unit may stand for a component or smaller entities (package, class, method)
				\item tests created by the developers
				\item automation is common (e.g., JUnit)
		\end{itemize}}
	\end{definition}
	\begin{definition}{2. Integration Testing \deutsch{Integrationstest}}
		\uncover<3->{\begin{itemize}
				\item some components are integrated (e.g., into subsystems) and tested together
				\item detects inconsistencies in interfaces and communication between components
				\item top-down vs bottom-up integration
		\end{itemize}}
	\end{definition}
	\nextcolumn
	\begin{definition}{3. System Testing \deutsch{Systemtest}}
		\uncover<4->{\begin{itemize}
				\item all components are integrated to the complete system
				\item detects further inconsistencies and unanticipated interactions
				\item system is tested against system requirements
		\end{itemize}}
	\end{definition}
	\begin{definition}{4. Acceptance Testing \deutsch{Abnahmetest}}
		\uncover<5->{\begin{itemize}
				\item final stage in the testing process before accepted for operational use
				\item system is tested against user requirements and with real data
				\item performed by (potential) customer
		\end{itemize}}
	\end{definition}
\end{fancycolumns}
}

% introduced: se1-versioncontrol
% reused: se2-configmanagement
\newcommand{\slideVersionControl}{
\begin{fancycolumns}[animation=none]
	\begin{note}{Goals of Version Control}
		\begin{itemize}
			\item collaborative work: synchronize files and folders with other users
			\item compare own version with other versions
			\item merge files edited by several users
			\item history: access old versions, log changes
		\end{itemize}
	\end{note} \pause
	\begin{definition}{Version Control Systems}
		\begin{itemize}
			\item local only: SCCS (1972), RCS (1982), \ldots
			\item centralized: CVS (1986), SVN (2000), \ldots
			\item distributed: \emph{git} (2005), mercurial (2005), \ldots
		\end{itemize}
		\vspace{2mm}
		invented for software, useful for most files %(cf.\ agile methods)
	\end{definition} \pause
	\begin{example}{}
		Alternatives? Why not use locks (cf.\ databases)?
	\end{example} \pause
	\nextcolumn
	\picDark[width=\linewidth]{versioncontrol/trends-git} \pause
	\vspace{-4mm}
	\begin{example}{Personal Experience}
		\begin{itemize}
			\item 2004: started working with CVS in industry
			\item 2007: initiated research protoype FeatureIDE with SVN
			\item 2009: published FeatureIDE open source (history neglected)
			\item 2014: migrated FeatureIDE's code to git
		\end{itemize}
	\end{example}
\end{fancycolumns}
}

% introduced: se1-versioncontrol
% reused: se2-configmanagement
\newcommand{\slideDistributedVC}{
\vspace{-2mm}
\begin{fancycolumns}
	\peertopeer{computer}{tonne}{computer}{computer}{GIT Server}{GIT Server/Client X}{GIT Server/Client Y}{
		\draw[kante] (client1) to (server);
		\draw[kante] (server) to (client1);
		\draw[kante] (client2) to (server);
		\draw[kante] (server) to (client2);
		\draw[innerekante] (client1.80) to (client1.100);
		\draw[innerekante] (client1.-100) to (client1.-80);
		\draw[innerekante] (client2.80) to (client2.100);
		\draw[innerekante] (client2.-100) to (client2.-80);
	}
	\begin{definition}{Distributed Version Control Systems}
		\begin{itemize}
			\item peer-to-peer architecture
			\item clients use repository clone to synchronize
			\item version history available on each client
			%\item comparing client version with server version / history of repository clone
			%\item merging files edited by several clients using an updated version of the repository clone
		\end{itemize}
	\end{definition}
	\nextcolumn
	\only<2|handout:0>{\peertopeer{computer}{tonne}{computer}{computer}{Remote Repository: Revision 1-10}{~}{~}{
	}}
	\only<3|handout:0>{\peertopeer{computer}{tonne}{computer}{computer}{Remote Repository: Revision 1-10}{Repository Clone: 1-10}{Repository Clone: 1-10}{
			\draw[kante] (server) to (client1);
			\draw[kante] (server) to (client2);
	}}
	\only<4|handout:0>{\peertopeer{computer}{tonne}{computer}{computer}{Remote Repository: Revision 1-10}{Working Copy: 10\\Repository Clone: 1-10}{Working Copy: 5\\Repository Clone: 1-10}{
			\draw[innerekante] (client1.80) to (client1.100);
			\draw[innerekante] (client2.80) to (client2.100);
	}}
	\only<5|handout:0>{\peertopeer{computer}{tonne}{computer}{computer}{Remote Repository: Revision 1-10}{Working Copy: 10 (changed)\\Repository Clone: 1-10}{Working Copy: 5\\Repository Clone: 1-10}{
	}}
	\only<6|handout:0>{\peertopeer{computer}{tonne}{computer}{computer}{Remote Repository: Revision 1-10}{Working Copy: 11\\Repository Clone: 1-11}{Working Copy: 5\\Repository Clone: 1-10}{
			\draw[innerekante] (client1.-100) to (client1.-80);
	}}
	\only<7|handout:1>{\peertopeer{computer}{tonne}{computer}{computer}{Remote Repository: Revision 1-11}{Working Copy: 11\\Repository Clone: 1-11}{Working Copy: 5\\Repository Clone: 1-10}{
			\draw[kante] (client1) to (server);
	}}
	\only<8|handout:0>{\peertopeer{computer}{tonne}{computer}{computer}{Remote Repository: Revision 1-11}{Working Copy: 11\\Repository Clone: 1-11}{Working Copy: 5\\Repository Clone: 1-11}{
			\draw[kante] (server) to (client2);
	}}
	\only<9|handout:0>{\peertopeer{computer}{tonne}{computer}{computer}{Remote Repository: Revision 1-11}{Working Copy: 11\\Repository Clone: 1-11}{Working Copy: 11\\Repository Clone: 1-11}{
			\draw[innerekante] (client2.80) to (client2.100);
	}}
	\begin{note}{Distributed Version Control Systems}
		\begin{itemize}
			\item main repository on one or several servers
			\item clone of the repository on each client 
			\item new revision on every change on the clients
			%\item Revisions used to undo changes and to merge files
		\end{itemize}
	\end{note}
\end{fancycolumns}
}

% introduced: se1-versioncontrol
% reused: se2-configmanagement
\newcommand{\slideConfigManagement}{
	\begin{fancycolumns}
		\begin{definition}{Configuration Management}
			\mycite{\emph{Configuration management} is concerned with the policies, processes, and tools for managing changing software systems.} % TODO fix source Aiello and Sachs 2011
		\end{definition}
		\begin{definition}{Four Activities in Configuration Management}
			\begin{itemize}
				\item version control / management: keeping track of multiple versions, enabling simultaneous changes
				\item system building: collecting, compiling, and linking components into executable systems
				\item change management: tracking change requests and planning if/when realized
				\item release management: preparing new and managing old releases
			\end{itemize}
		\end{definition}
		\nextcolumn
		\begin{definition}{Development Stages \deutsch{Entwicklungsphasen}}
			\begin{itemize}
				\item \emph{development phase}: adding new functionality
				\item \emph{system testing phase}: internal release, bug fixes, performance improvements, security fixes, but no new functionality
				\item \emph{release phase}: bug reports and feature requests by users or customers
				\item often several versions co-exist at different stages
			\end{itemize}
		\end{definition}
	\end{fancycolumns}	
}

% introduced: se1-versioncontrol
% reused: se2-configmanagement
\newcommand{\slideCloneFetch}{
\begin{frame}{Clone and Fetch}
	\begin{fancycolumns}[animation=none,T]
		\only<1|handout:0>{\peertopeer{ordner1}{tonne}{empty}{empty}{Repository: Revision 1-10}{~}{~}{}}%
		\only<2|handout:1>{\peertopeer{ordner1}{tonne}{empty}{empty}{Repository: Revision 1-10}{Repository Clone: 1-10}{Repository Clone: 1-10}{
				\draw[kante] (server) to node[swap] {Clone Repository} (client1);
				\draw[kante, bend right, <-] (client2) to node[swap] {Clone Repository} (server);
		}}%
		\only<3|handout:0>{\peertopeer{ordner1}{tonne}{ordner1}{ordner1}{Repository: Revision 1-10}{Working Copy: 10\\Repository Clone: 1-10}{Working Copy: 5\\Repository Clone: 1-10}{
				\draw [innerekante, <-, in = 135, out = 90]  ([xshift = 20pt, yshift = 5pt]client1.north west) to node[swap, above] {Working Copy: 10} ([xshift = -10pt] client1.north east);
				
				\draw [innerekante, <-, in = 135, out = 90]  ([xshift = 20pt, yshift = 5pt]client2.north west) to node[swap, above] {Working Copy: 5} ([xshift = -10pt] client2.north east);
		}}%
		\only<4-|handout:0>{\peertopeer{ordner1}{tonne}{ordner1}{ordner1}{Repository: Revision 1-10}{Working Copy: 10\\Repository Clone: 1-10}{Working Copy: 5\\Repository Clone: 1-10}{
		}}%
		\only<1-|handout:1>{
			\begin{note}{Clone}
				\begin{itemize}
					\setlength{\itemsep}{.1em}
					\item create a local repository clone 
					\item use local repository to create a working copy
					\item specify folders for the clone
					\item specify revision or head (latest revision)
				\end{itemize}
			\end{note}
		}
		\nextcolumn
		\only<5|handout:0>{\peertopeer{ordner2}{tonne}{ordner2}{ordner1}{Repository: Revision 1-11}{Working Copy: 11\\Repository Clone: 1-11}{Working Copy: 5\\Repository Clone: 1-10}{}}	
		\only<6|handout:1>{\peertopeer{ordner2}{tonne}{ordner2}{ordner1}{Repository: Revision 1-11}{Working Copy: 11\\Repository Clone: 1-11}{Working Copy: 5\\Repository Clone: 1-10}{
				\draw[kante, ->] (server) to node[swap, above right] {Fetch Repository} (client2);
		}}
		\only<7|handout:0>{\peertopeer{ordner2}{tonne}{ordner2}{ordner1}{Repository: Revision 1-11}{Working Copy: 11\\Repository Clone: 1-11}{Working Copy: 5\\Repository Clone: 1-\emph{11}}{
		}}
		\only<5-7|handout:1>{
			\begin{note}{Fetch}
				\begin{itemize}
					\setlength{\itemsep}{.1em}
					\item download the remote repository
					\item note: working copy remains unchanged
				\end{itemize}
			\end{note}
		}
	\end{fancycolumns}
\end{frame}
}

% introduced: se1-versioncontrol
% reused: se2-configmanagement
\newcommand{\slideCommitPush}{
\begin{frame}{Commit and Push}
	\begin{fancycolumns}[animation=none,T]
		\only<1|handout:0>{\peertopeer{ordner1}{tonne}{ordner2}{ordner1}{Repository: Revision 1-10}{Working Copy: 10 \emph{(changed)}\\Repository Clone: 1-10}{Working Copy: 5\\Repository Clone: 1-10}{
		}}
		\only<2|handout:1>{\peertopeer{ordner1}{tonne}{ordner2}{ordner1}{Repository: Revision 1-10}{Working Copy: 10 (changed)\\Repository Clone: 1-10}{Working Copy: 5\\Repository Clone: 1-10}{
				\draw [innerekante, ->, in = 135, out = 90]  ([xshift = 20pt, yshift = 5pt]client1.north west) to node[swap, above] {Commit "Change message"} ([xshift = -10pt] client1.north east);
		}}
		\only<3-|handout:0>{\peertopeer{ordner1}{tonne}{ordner2}{ordner1}{Repository: Revision 1-10}{Working Copy: \emph{11}\\Repository Clone: 1-\emph{11}}{Working Copy: 5\\Repository Clone: 1-10}{
		}}
		
		\only<1-|handout:1>{
			\begin{note}{Commit}
				\begin{itemize}
					\setlength{\itemsep}{.1em}
					\item commit changes to local repository
					\item commits are atomic (all or nothing)
					\item specify changed folders and files
					\item mandatory message describing your change
				\end{itemize}
			\end{note}
		}
		\nextcolumn
		\only<4|handout:0>{\peertopeer{ordner1}{tonne}{ordner2}{ordner1}{Repository: Revision 1-10}{Repository Clone: 1-\emph{11}}{Working Copy: 5}{}}%
		\only<5|handout:1>{\peertopeer{ordner1}{tonne}{ordner2}{ordner1}{Repository: Revision 1-10}{Repository Clone: 1-11}{Working Copy: 5}{
				\draw[kante, <-] (server) to node[swap] {Push Repository} (client1);
		}}%
		\only<6|handout:0>{\peertopeer{ordner2}{tonne}{ordner2}{ordner1}{Repository: Revision 1-\emph{11}}{Repository Clone: 1-11}{Working Copy: 5}{
		}}%
		\only<4-6|handout:1>{
			\begin{note}{Push}		
				\begin{itemize}
					\setlength{\itemsep}{.1em}
					\item push local repository to server
					\item condition: local repository is up to date (might require previous pull)
				\end{itemize}
			\end{note}
		}
	\end{fancycolumns}
\end{frame}
}	
% introduced: se1-versioncontrol
% reused: se2-configmanagement
\newcommand{\slideMergePull}{
\begin{frame}{Merge and Pull}
	\vspace{-1.5mm}
	\begin{fancycolumns}[animation=none,T]
		\only<1|handout:1>{\peertopeer{ordner2}{tonne}{ordner2}{ordner1}{Remote Repository: Revision 1-11}{Working Copy: 11\\Repository Clone: 1-11}{Working Copy: 5\\Repository Clone: 1-11}{
				\draw [innerekante, <-, in = 135, out = 90] ([xshift = 20pt, yshift = 5pt]client2.north west) to node[swap, above] {Merge to 11} ([xshift = -10pt]client2.north east);
		}}%
		\only<2-|handout:0>{\peertopeer{ordner2}{tonne}{ordner2}{ordner1}{Remote Repository: Revision 1-11}{Working Copy: 11\\Repository Clone: 1-11}{Working Copy: \emph{11}\\Repositoy Clone: 1-11}{
		}}%
		\only<1-|handout:1>{
			\begin{note}{Merge}
				\begin{itemize}
					\item update working copy with local repository
					\item integrate new commits from other branch
					\item fast forward / automatic merge / manual merge
				\end{itemize}
			\end{note}
		}
		\nextcolumn
		\only<3|handout:0>{\peertopeer{ordner2}{tonne}{ordner2}{ordner1}{Remote Repository: Revision 1-\emph{11}}{Working Copy: 11\\Repository Clone: 1-11}{Working Copy: 5\\Repository Clone: 1-10}{
				\draw[kante, ->] (server) to node[swap] {Fetch Repository} (client2);
		}}
		\only<4|handout:0>{\peertopeer{ordner2}{tonne}{ordner2}{ordner1}{Remote Repository: Revision 1-11}{Working Copy: 11\\Repository Clone: 1-11}{Working Copy: 5\\Repository Clone: 1-\emph{11}}{
		}}
		\only<5|handout:0>{\peertopeer{ordner2}{tonne}{ordner2}{ordner1}{Remote Repository: Revision 1-11}{Working Copy: 11\\Repository Clone: 1-11}{Working Copy: 5\\Repository Clone: 1-11}{
				\draw [innerekante, <-, in = 135, out = 90] ([xshift = 20pt, yshift = 5pt]client2.north west) to node[swap, above] {Merge to 11} ([xshift = -10pt]client2.north east);
		}}
		\only<6|handout:0>{\peertopeer{ordner2}{tonne}{ordner2}{ordner1}{Remote Repository: Revision 1-11}{Working Copy: 11\\Repository Clone: 1-11}{Working Copy: \emph{11}\\Repository Clone: 1-11}{
		}}	
		\only<-0|handout:1>{\peertopeer{ordner2}{tonne}{ordner2}{ordner1}{Remote Repository: Revision 1-11}{Working Copy: 11\\Repository Clone: 1-11}{Working Copy: \emph{11}\\Repository Clone: 1-\emph{11}}{
				\draw[kante, ->] (server) to node[swap] {Fetch Repository} (client2);
				\draw [innerekante, <-, in = 135, out = 90] ([xshift = 20pt, yshift = 5pt]client2.north west) to node[swap, above] {Merge to 11} ([xshift = -10pt]client2.north east);
		}}
		\only<3-6|handout:1>{
			\begin{note}{Pull}
				\begin{itemize}
					\item pull = fetch + merge
					\item download and merge in one step!
				\end{itemize}
			\end{note}
		}
	\end{fancycolumns}
\end{frame}
}
% introduced: se1-versioncontrol
% reused: se2-configmanagement
\newcommand{\slideIgnore}{
\begin{frame}
	\frametitle{Ignore}
	\begin{fancycolumns}[widths={55},animation=none]
		\only<1|handout:0>{\peertopeer{ordner1}{tonne}{ordner2}{ordner1}{Remote Repository: Revision 1-10}{Working Copy: 10 \emph{(changed)}\\Repository Clone: 1-10}{Working Copy: 5\\Repository Clone: 1-10}{
		}}
		\only<2|handout:1>{\peertopeer{ordner1}{tonne}{ordner2}{ordner1}{Remote Repository: Revision 1-10}{Working Copy: 10 (changed)\\Repository Clone: 1-10}{Working Copy: 5\\Repository Clone: 1-10}{
				\draw [innerekante, ->, in = 135, out = 90] ([xshift = 20pt, yshift = 5pt]client1.north west) to node[swap, above] {Ignore Foo.class + Commit} ([xshift=-10pt]client1.north east);
		}}
		\only<3|handout:0>{\peertopeer{ordner1}{tonne}{ordner2}{ordner1}{Remote Repository: Revision 1-10}{Working Copy \emph{11}\\Repository Clone: 1-\emph{11}}{Working Copy: Revision 5\\Repository Clone: 1-10}{
		}}
		\only<4|handout:0>{\peertopeer{ordner1}{tonne}{ordner2}{ordner1}{Remote Repository: Revision 1-10}{Working Copy: 11\\Repository Clone: 1-11}{Working Copy: 5\\Repository Clone: 1-10}{
				\draw[kante, <-] (server) to node[swap] {Push repository} (client1);
		}}
		\only<5|handout:0>{\peertopeer{ordner1}{tonne}{ordner2}{ordner1}{Remote Repository: Revision 1-\emph{11}}{Working Copy: 11\\Repository Clone: 1-11}{Working Copy: 5\\Repository Clone: 1-10}{
		}}		
		\begin{note}{.gitignore}
			\begin{itemize}
				\item resources will be ignored on commit
				\item user-specific files or derived/generated resources
				\item specify files and folders (e.g., *.log files)
			\end{itemize}
		\end{note}
	\end{fancycolumns}
\end{frame}
}

% introduced: se1-versioncontrol
% reused: se2-configmanagement
\newcommand{\frameSystemBuilding}{
\begin{fancycolumns}
	\begin{definition}{System Building}
		\mycite{\emph{System building} is the process of creating a complete, executable system by compiling and linking the system components, external libraries, configuration files, and other information.}
	\end{definition}
	\begin{definition}{Building Involves Three Platforms}
		\begin{itemize}
			\item \emph{development system}: compilers and editors used on the developer's system to test prior to commit
			\item \emph{build server}: server to build and distribute executable versions, triggered by commits or schedule (i.e., nightly builds)
			\item \emph{target environment}: intended platform for executable system (e.g., ECU in a car)
		\end{itemize}
	\end{definition}
	\nextcolumn
	\begin{definition}{Tooling for Building and System Integration}
		\begin{itemize}
			\item build script generation: identify dependent components, automated generation or tool support for creation and editing
			\item version control system integration: checkout required versions of components
			\item minimal recompilation: determine which parts need to be recompiled
			\item executable system creation: compilation and linking
			\item test automation: run automated tests (e.g., unit tests)
			\item reporting: reports about success or failure of builds and tests
			\item documentation generation: release notes, help pages
		\end{itemize}
	\end{definition}
\end{fancycolumns}
}


% introduced: se1-versioncontrol
% reused: se2-configmanagement
\newcommand{\frameContinousIntegration}{
	\begin{fancycolumns}
		\begin{note}{Continuous Integration}
			\mycite{Agile methods recommend that very frequent system builds should be carried out, with automated testing used to discover software problems. Frequent builds are part of a process of continuous integration [...].}
		\end{note}
		\begin{definition}{Continuous Integration Tools}
			\begin{itemize}
				\item \href{https://en.wikipedia.org/wiki/CruiseControl}{CruiseControl} (2001--2010) --- first open-source tool
				\item \href{https://en.wikipedia.org/wiki/TeamCity}{TeamCity} (2006--)
				\item \href{https://en.wikipedia.org/wiki/Hudson_(software)}{Hudson} (2008--2017)
				\item \href{https://en.wikipedia.org/wiki/Jenkins_(software)}{Jenkins} (2011--) --- fork of Hudson
				\item \href{https://en.wikipedia.org/wiki/Travis_CI}{Travis CI} (2011--) --- made in Germany
				\item \href{https://en.wikipedia.org/wiki/GitLab}{GitLab} (2014--)
				% TODO add more tools? go continuous delivery, codehsip, circleci, codefresh, bamboo
			\end{itemize}
		\end{definition}
		\nextcolumn
		\begin{definition}{Steps in Continuous Integration}
			\begin{itemize}
				\item clone/fetch from version control
				\item if feasible: build and run automated tests, if it fails others are responsible
				\item apply changes
				\item build and run automated tests locally, if it fails continue editing
				\item if local tests pass, commit to feature branch in version control
				\item commit triggers build server, if it fails continue editing
				\item if tests pass (and code review approves changes), merge branch into main development branch
			\end{itemize}
		\end{definition}
	\end{fancycolumns}
}

% introduced: se1-versioncontrol
% reused: se2-configmanagement
\newcommand{\frameDevOps}{
\begin{fancycolumns}
	\begin{note}{Motivation}
		\begin{itemize}
			\item if software fails: 
			\begin{itemize}
				\item programmers blame administrators for misconfiguration
				\item administrators blame programmers for erroneous software
			\end{itemize}
			\item programmers want frequent updates
			\item administrators follow the slogan: \mycite{never change a running system}
			\item customers and users want a single responsibility
			\item shorter and shorter update cycles
		\end{itemize}
	\end{note}
	\nextcolumn
	\begin{definition}{DevOps}
		\begin{itemize}
			\item promoted in agile development
			\item \emph{Dev}: development by programmers
			\item \emph{Ops}: operation \deutsch{Betrieb} by administrators
			\item \emph{DevOps}: teams that are responsible for both, development and operations
			\item goal: avoid blaming each other by shared responsibility
		\end{itemize}
	\end{definition}
\end{fancycolumns}
}

% introduced: se2-introduction
% reused: se2-summary
\newcommand{\frameCourseQuestions}{
	\small
	\begin{fancycolumns}[columns=3,t]
		\begin{definition}{Part A: How to \emph{change} software?}
			\begin{itemize}
				\item Evolution: How does software change over time and why?
				\item Maintenance: How is software operated and maintained?
				\item Design Patterns: How to better support future changes?
				\item (Lego Scrum: How to develop agile and in teams?)
			\end{itemize}
		\end{definition}
		\nextcolumn
		\begin{definition}{Part B: How to \emph{verify} software?}
			\begin{itemize}
				\item Compilation: How do compilers work and what is their role in software engineering?
				\item Static Analysis: How can we detect faults without executing the software?
				\item Dynamic Analysis: How can we detect faults by executing the software?
				\item Design by Contract: How to specify the intended behavior of the software?
				\item Configuration Management: How to manage and verify versions of software?
			\end{itemize}
		\end{definition}
		\nextcolumn
		\begin{definition}{Part C: How to \emph{reuse} software?}
			\begin{itemize}
				\item Open-Source Software: When are we allowed to reuse software?
				\item Software Product Lines: How can we reuse software systematically?
				\item Automotive Software: How is reuse is performed for automotive software?
				\item Guest Lectures: How is software engineering done in practice and envisioned in research?
			\end{itemize}
		\end{definition}
	\end{fancycolumns}
}

% introduced: se2-evolution
% reused: se2-summary
\newcommand{\frameLehmansLaws}{
	\begin{fancycolumns}
		\begin{definition}{\insertsubsection}
			\begin{itemize}
				\item Continuing Change (1974)
				\item Increasing Complexity (1974)
				\item Conservation of Familiarity (1980)
				\item Continuing Growth (1980)
				\item Declining Quality (1996)
				\item Feedback System (1996)
				\item laws not discussed in this course:\\Self Regulation,\\Conservation of Organizational Stability,\\\ldots
			\end{itemize}
		\end{definition}
		\nextcolumn
		\begin{note}{Essence of the Laws}
			\begin{itemize}
				\item software that is used will be modified
				\item when modified, its complexity will increase (unless one does actively work against it)
			\end{itemize}
		\end{note}
		\begin{example}{Consequences}
			\begin{itemize}
				\item functional changes are inevitable
				\item changes are not necessarily a consequence of errors (e.g., in requirements engineering or programming)
				\item there are limits to what a development team can achieve (cf.\ Continuing Growth)
			\end{itemize}
		\end{example}
	\end{fancycolumns}
}


% UNIVERSITY-SPECIFIC ADJUSTMENTS
% usage: \def\university{ulm}

\ifdefined\university
\else
\def\university{}
\fi

\newcommand{\ifuniversity}[3][]{\ifthenelse{\equal{\university}{#2}}{#3}{#1}}
\newcommand{\unlessuniversity}[3][]{\ifuniversity[#3]{#2}{#1}}
\newcommand{\inputuniversity}[1]{\IfFileExists{#1-\university}{\input{#1-\university}}{}}

% TITLE
\title{Introduction to Software Engineering} % default title for the course
\ifuniversity{tubs}{\title[SE1]{Software Engineering 1}}

% LECTURERS AND TUTORS
\makeatletter
\let\author@old\author
\renewcommand{\author}[1]{
	\ifuniversity{tubs}{\author@old{Thomas Thüm, Christopher Rau}}
	\ifuniversity{recording}{\author@old{#1}}
	\ifuniversity{anonymous}{\author@old{Anonymous Authors}}
	\ifuniversity{}{\author@old{#1}}
}
\makeatother

% CONTENT OVERVIEW
\newcommand{\lectureseriesoverview}[1][]{ % shown at the start and end of each lecture
	\setlength{\leftmargini}{3.5ex}
	\begin{fancycolumns}[t,columns=3,animation=none]
		\mydefinition{\emphifequal{#1}{I}{Part A: Development of \mbox{\emph{Correct} Software}}}{
			\begin{enumerate}
				\item {\emphifequal{#1}{1}{Introduction}}
				\item {\emphifequal{#1}{2}{Implementation}}
				\item {\emphifequal{#1}{3}{Testing}}
				\item {\emphifequal{#1}{4}{Maintenance}}
%				\item {\emphifequal{#1}{4}{Software Changes}}
%				\item {\emphifequal{#1}{5}{Version Control}}
			\end{enumerate}
		}
		\nextcolumn
		\mydefinition{\emphifequal{#1}{II}{Part B: Development in \emph{Schedule} \& \emph{Budget}}}{
			\begin{enumerate}
				\setcounter{enumi}{4}
				\item {\emphifequal{#1}{5}{Project Management}}
				\item {\emphifequal{#1}{6}{Development Process}}
			\end{enumerate}
		}
		\nextcolumn
		\mydefinition{\emphifequal{#1}{III}{Part C: Development of \mbox{\emph{Needed} Software}}}{
			\begin{enumerate}
				\setcounter{enumi}{6}
				\item {\emphifequal{#1}{7}{Requirements}}
				\item {\emphifequal{#1}{8}{System Modeling}}
				\item {\emphifequal{#1}{9}{Software Architecture}}
				\item {\emphifequal{#1}{10}{Software Design}}
				\item {\emphifequal{#1}{11}{Software Reuse}}
				\item {\emphifequal{#1}{12}{Guest Lecture}}
			\end{enumerate}
		}
	\end{fancycolumns}
}
\newcommand{\skipPartA}{\addtocounter{lecture}{4}}
\newcommand{\skipPartB}{\addtocounter{lecture}{2}}

% COLOR SCHEME, LOGOS, PICTURES
\newcommand*{\lectureqr}[2][]{} % as default there is no qr code

\ifuniversity{tubs}{
	\setpaths{
		{../pics/tubs/}
		{../pics/tubs/2024wt/}
	}
	
	\renewcommand*{\lectureqr}[2][width=.33\linewidth]{
		\begin{flushright}
			\picDark[#1]{qr#2}
		\end{flushright}
	}
	
	\usepackage{tubscolors}
	% redefine colors used in fancybeamer
	% Orange, OrangeLight (Yellow), OrangeDark (tubsRed)
	% Green, GreenLight, GreenDark
	% Blue, BlueLight, BlueDark
	% Violet, VioletLight, VioletDark
	% tubsGray (tubsGray60), tubsLightGray (tubsGray20)
	% tubsSecondaryLight, tubsSecondaryMedium, tubsSecondary, tubsSecondaryDark
	% x20, x40, x60, x80, x100

	\AtEndPreamble{
		\colorlet{red}{tubsRed}
		\colorlet{accent}{background}
		\colorlet{accenttwo}{red}
		\setbeamercolor{logobox}{bg=background}
		\setbeamercolor{myfooter}{fg=foreground,bg=accent}
		\setbeamercolor{subtitlebox}{fg=foreground,bg=accent}
		\fancylogos{tubs,isf}
		\ifdarkmode
			\colorlet{red}{red!90!white}
			\colorlet{green}{green!90!white}
			\colorlet{blue}{blue!90!white}
			\colorlet{orange}{orange!90!white}
			\colorlet{accenttwo}{accenttwo!25!darkgray}
			\setbeamercolor{mypagenumber}{fg=foreground,bg=accenttwo}
			\setbeamercolor{titlebox}{fg=foreground,bg=accenttwo}
			\fancylogos{tubs-black,isf}
		\fi

		% hack to force use of changed colors
		\UpdateBoxColor{definition}{orange}
		\UpdateBoxColor{example}{blue}
		\UpdateBoxColor{note}{red}
	}

	\setpicture[150]{campus/04-25-iz1} % default picture
}

\ifuniversity{anonymous}{
	\fancylogos{}
	\renewcommand{\pic}[2][]{\includegraphics[#1]{#2}}
	\renewcommand{\deutsch}[1]{}
	\renewcommand{\deutschertitel}[1]{}
}

\ifuniversity{recording}{
	\renewcommand{\deutsch}[1]{} % no German in English recordings
	\mode<beamer>{\renewcommand{\pic}[2][]{\includegraphics[#1]{#2}}} % avoid annoying tool tips during the recording
	%	\setpicture[550]{jun22-clouds3} % default picture
}{
	\AtEndPreamble{
		\hypersetup{colorlinks=true, linkcolor=red, urlcolor=red} % emphasizing links, but not when recording
	}
}
 % this file is specific to the respective lecture

% HOPEFULLY NON-PERMANENT HACKS

\usepackage{pgf, tikz}
\usetikzlibrary{arrows, arrows.meta, positioning, shapes, calc, decorations.text, shadows.blur, matrix, fit, mindmap, backgrounds}
\pgfdeclarelayer{foreground}
\pgfsetlayers{background,main,foreground}
\tikzset{ >=stealth'} % define standard arrow tip
\usepackage{tikz-qtree}

\newcommand{\handoutonly}[1]{}
\newcommand{\slideonly}[1]{#1}
\mode<handout>{
	\renewcommand{\handoutonly}[1]{#1}
	\renewcommand{\slideonly}[1]{}
}
\newcommand{\prehandout}[1]{#1}
\newcommand{\posthandout}[1]{}
\newcommand{\onlyhandout}[1]{\ifdefined\ishandout#1\fi}

\renewcommand{\prehandout}[1]{}
\renewcommand{\posthandout}[1]{#1}

% Keys to support piece-wise uncovering of elements in TikZ pictures:
% \node[visible on=<2->](foo){Foo}
% \node[visible on=<{2,4}>](bar){Bar}   % put braces around comma expressions
%
% Internally works by setting opacity=0 when invisible, which has the 
% adavantage (compared to \node<2->(foo){Foo} that the node is always there, hence
% always consumes space plus that coordinate (foo) is always available.
%
% The actual command that implements the invisibility can be overriden
% by altering the style invisible. For instance \tikzsset{invisible/.style={opacity=0.2}}
% would dim the "invisible" parts. Alternatively, the color might be set to white, if the
% output driver does not support transparencies (e.g., PS) 
%
\tikzset{
	invisible/.style={opacity=0},
	visible on/.style={alt={#1{}{invisible}}},
	alt/.code args={<#1>#2#3}{%
		\alt<#1>{\pgfkeysalso{#2}}{\pgfkeysalso{#3}} % \pgfkeysalso doesn't change the path
	},
}

\usepackage{tabularx}



% QR Code handling
\usepackage{fancyqr}
\makeatletter
\let\oldfancyqr\fancyqr
\renewcommand{\fancyqr}[2]{\renewcommand{\@minipagerestore}{\setlength{\parskip}{0pt}}\oldfancyqr[#1]{#2}} % TODO: enable []
\makeatother

\newcommand{\studipqr}{
	\fancyqr{color=black,height=3.5cm}{\StudIPLink}
}

